\documentclass[12pt,a4paper]{report}
\usepackage[utf8]{inputenc}
\usepackage[T1]{fontenc}
\usepackage{lmodern}
\usepackage{geometry}
\geometry{margin=1in}
\usepackage{hyperref}
\usepackage{listings}
\usepackage{xcolor}
\usepackage{amsmath}
\usepackage{titlesec}
\usepackage{fancyhdr}
\usepackage{booktabs}
\usepackage{array}

\titleformat{\chapter}[display]
  {\normalfont\huge\bfseries}{\chaptertitlename\ \thechapter}{20pt}{\Huge}
\titlespacing*{\chapter}{0pt}{50pt}{40pt}

\lstset{
  basicstyle=\ttfamily\footnotesize,
  breaklines=true,
  frame=single,
  backgroundcolor=\color{gray!10},
  keywordstyle=\color{blue},
  commentstyle=\color{green!50!black},
  stringstyle=\color{red}
}

\pagestyle{fancy}
\fancyhf{}
\rhead{Abyss OS Theory Guide}
\lhead{\leftmark}
\rfoot{Page \thepage}

\title{\textbf{Abyss OS: Complete Architecture and Theory Guide}\\[1em]
\large A Deep-Dive into Bare-Metal x86 Operating System Design}
\author{Built by After-Dark \& Antigravity}
\date{\today}

\begin{document}
\maketitle
\tableofcontents
\newpage

\chapter{Toolchain \& Build System}

\section{The Cross-Compilation Problem}

When you compile a C++ program on a standard Linux machine, the compiler performs
a series of hidden assumptions. It links against \texttt{glibc} (the GNU C standard
library), injects startup code from \texttt{crt0.o} that initializes the C runtime
environment, and produces an ELF binary that depends on a Linux kernel being present
to load and execute it.

None of these assumptions hold for a bare-metal kernel. There is no Linux beneath us,
no \texttt{glibc} to link against, and no OS loader to parse the ELF binary. If we
attempted to run a standard Linux-compiled binary directly on bare hardware, the CPU
would execute the startup code, attempt a system call to the kernel to set up memory
regions, receive no response, and immediately generate a General Protection Fault.

The solution is a \textbf{cross-compiler}: a compiler that \textit{runs} on the host
machine (Linux x86-64) but \textit{targets} a completely different, standalone
environment (\texttt{i686-elf}). The \texttt{elf} in the target triple means the
output format is ELF32, but with no OS ABI, no linked libraries, and no startup
overhead whatsoever.

We compiled GCC and GNU Binutils from source with the following critical flags:
\begin{itemize}
    \item \textbf{\texttt{--target=i686-elf}}: Targets a 32-bit x86 bare-metal environment.
    \item \textbf{\texttt{--without-headers}}: Prevents GCC from using any OS header files.
    \item \textbf{\texttt{--disable-nls}}: Removes internationalization support (unnecessary overhead).
\end{itemize}

\section{Compiler Flags for Bare-Metal}

The flags passed to \texttt{i686-elf-g++} enforce the strict discipline of
bare-metal programming:

\begin{itemize}
    \item \textbf{\texttt{-ffreestanding}}: Informs the compiler that the standard
    library does not exist. It will not generate implicit calls to library functions,
    even for seemingly simple operations like zero-initializing a large array.
    \item \textbf{\texttt{-fno-exceptions}}: C++ exception handling (\texttt{try/catch})
    requires the compiler to generate \textit{unwind tables} and link against the
    \texttt{libgcc\_eh} library for stack unwinding. Since our kernel lacks that library,
    we disable this feature entirely.
    \item \textbf{\texttt{-fno-rtti}}: Run-Time Type Information (the mechanism behind
    \texttt{dynamic\_cast} and \texttt{typeid}) requires a hidden table of type metadata
    to be embedded in the binary and queried at runtime. This is unnecessary overhead
    for a kernel and is disabled.
    \item \textbf{\texttt{-O2}}: Instructs the compiler to apply a standard suite of
    optimizations --- loop unrolling, constant folding, dead-code elimination ---
    to produce faster machine code without making debugging prohibitively difficult.
\end{itemize}

\section{The Build System and Linker Script}

A \textbf{linker script} (\texttt{linker.ld}) controls the physical memory layout of
the final binary. When the CPU boots, it enters Real Mode at physical address
\texttt{0x7C00} (where the bootloader lives). The GRUB bootloader then loads our
kernel binary starting at physical address \texttt{0x100000} (the 1MB mark), which is
the conventional load address for 32-bit kernels.

Our linker script enforces this layout:
\begin{lstlisting}[title=linker.ld]
ENTRY(_start)

SECTIONS {
    . = 1M;           /* Set the load address to exactly 1 Megabyte */

    .text BLOCK(4K) : ALIGN(4K) {
        *(.multiboot)  /* Multiboot header MUST come first */
        *(.text)       /* All executable code sections */
    }
    .data BLOCK(4K) : ALIGN(4K) { *(.data) }
    .bss  BLOCK(4K) : ALIGN(4K) { *(.bss)  }
}
\end{lstlisting}

The \texttt{BLOCK(4K)} directive aligns each section to a 4096-byte page boundary,
which is required by the Paging hardware that we implement in a later chapter.

\newpage

\chapter{Boot Sequence \& Display}

\section{The x86 Boot Process}

When an x86 computer is powered on, the CPU begins execution in \textbf{Real Mode},
a 16-bit operating environment with a 1MB address space that dates back to the Intel
8086 (1978). The CPU's instruction pointer is initialized to physical address
\texttt{0xFFFF0}, which points to the motherboard's \textbf{BIOS} (Basic Input/Output
System) firmware burned into ROM.

The BIOS performs a \textbf{Power-On Self-Test (POST)}: it enumerates and tests the
installed RAM, discovers peripheral devices, and initializes essential controllers.
It then searches connected storage devices for a bootable medium by reading the very
first 512-byte sector. If the last two bytes of that sector are the magic signature
\texttt{0x55 0xAA}, the BIOS copies those 512 bytes to address \texttt{0x7C00} and
transfers execution to it.

Abyss OS bypasses the 512-byte limitation entirely by relying on \textbf{GRUB} (GNU
GRand Unified Bootloader) as the first-stage loader. GRUB handles the transition
from 16-bit Real Mode to 32-bit Protected Mode, sets up the initial stack, and loads
our kernel binary from the ISO image.

\section{The Multiboot Specification}

GRUB uses the \textbf{Multiboot Specification} to identify a valid kernel binary.
Our kernel must embed a specific magic structure within its first 8KB, which GRUB
detects by scanning the binary. This structure contains:

\begin{itemize}
    \item \textbf{Magic Number} \texttt{0x1BADB002}: A fixed constant that GRUB scans for.
    \item \textbf{Flags}: A bitmask requesting features from GRUB (e.g., module loading,
    memory maps, video mode configuration).
    \item \textbf{Checksum}: A value such that \texttt{Magic + Flags + Checksum == 0}.
    This arithmetic check prevents GRUB from falsely identifying arbitrary data as a
    Multiboot header.
\end{itemize}

When GRUB successfully locates and validates the header, it switches the CPU into
32-bit Protected Mode, disables interrupts (\texttt{cli}), and calls our \texttt{\_start}
label. At this moment, GRUB passes a critical pointer in the \texttt{EBX} register:
the address of the \textbf{Multiboot Information Structure}, a rich data structure
containing the physical RAM map, framebuffer address, and loaded module locations.

\section{Protected Mode and Memory Segmentation}

\textbf{Protected Mode} is the 32-bit operating mode of the x86 architecture,
introduced by the Intel 80286. It expands the addressable memory space from 1MB
(20-bit address bus) to 4GB (32-bit address bus) and introduces hardware-enforced
memory protection via \textbf{Segmentation}.

In Protected Mode, memory is no longer addressed directly. Instead, all memory
accesses pass through a \textbf{Segment Register} (e.g., \texttt{CS} for code,
\texttt{DS} for data). These registers contain \textbf{selectors} --- indices into the
\textbf{Global Descriptor Table (GDT)}, which we explore in depth in Chapter 3.

\section{VGA Text Mode and Memory-Mapped I/O}

\textbf{VGA Text Mode} (specifically Mode 3) provides an 80-column $\times$ 25-row
character grid. The graphics hardware monitors a specific region of the CPU's physical
address space: \texttt{0xB8000}--\texttt{0xBFFFF}. This technique, where hardware
devices appear directly in the CPU's memory map, is called \textbf{Memory-Mapped I/O (MMIO)}.

Each character on screen is represented by two bytes in this region:
\begin{itemize}
    \item \textbf{Low byte}: The ASCII character code.
    \item \textbf{High byte}: The color attribute. Bits 0--3 are the foreground color;
    bits 4--6 are the background color; bit 7 enables blinking. For example,
    \texttt{0x0F} represents bright white (\texttt{0xF}) on black (\texttt{0x0}).
\end{itemize}

The VGA controller continuously \textit{scans} this buffer at the display refresh rate
(typically 60--70 Hz) and converts the ASCII codes to actual pixel glyphs using its
internal character generator ROM. The OS simply writes bytes to RAM; the hardware
handles all the actual rendering.

\begin{lstlisting}[language=C++, title=kernel/vga.cpp -- Core MMIO interaction]
volatile uint16_t* vga = (uint16_t*)0xB8000;

void vga_print_char(char c) {
    vga[cursor_pos] = (0x0F << 8) | c; // Color | ASCII
    cursor_pos++;
}
\end{lstlisting}

The \texttt{volatile} keyword is essential here. Without it, the C++ optimizer might
determine that the value written to memory is never \textit{read back} by the program
and eliminate the write entirely. \texttt{volatile} guarantees that every write is
committed to the physical address, preventing the optimizer from discarding MMIO writes.

\newpage

\chapter{CPU Architecture: GDT, IDT, PIC, and Exceptions}

\section{The Global Descriptor Table (GDT)}

\subsection{Segmentation and the Need for the GDT}

In Protected Mode, x86 uses a two-layer memory protection system: \textbf{Segmentation}
and \textbf{Paging}. Segmentation is the first layer. Every memory access the CPU makes
is relative to a \textit{segment}, which defines a contiguous region of memory with
specific access permissions.

The GDT is an array of \textbf{8-byte segment descriptors} that lives in RAM. Its
physical address is communicated to the CPU via the \texttt{lgdt} instruction, which
loads the address into the CPU's hidden \texttt{GDTR} register (Global Descriptor Table
Register).

\subsection{Descriptor Format}

Each 8-byte descriptor has a complex, non-contiguous bit layout for historical
compatibility reasons:

\begin{center}
\begin{tabular}{|c|c|c|c|}
\hline
\textbf{Bits} & \textbf{Field} & \textbf{Width} & \textbf{Purpose} \\
\hline
0--15   & Limit (low)  & 16 bits & Lower 16 bits of segment size \\
16--39  & Base (low)   & 24 bits & Lower 24 bits of base address \\
40--47  & Access byte  & 8 bits  & Privilege level, type, present flag \\
48--51  & Limit (high) & 4 bits  & Upper 4 bits of segment size \\
52--55  & Flags        & 4 bits  & 32-bit mode, granularity \\
56--63  & Base (high)  & 8 bits  & Upper 8 bits of base address \\
\hline
\end{tabular}
\end{center}

The \textbf{Access Byte} is the most important field. Its bit layout determines
the type and privilege of the segment:
\begin{itemize}
    \item \textbf{Bit 7 (Present)}: Must be 1 for valid descriptors.
    \item \textbf{Bits 5--6 (DPL)}: Descriptor Privilege Level (Ring 0--3).
    \item \textbf{Bit 4}: 1 for code/data segments, 0 for system segments.
    \item \textbf{Bit 3 (Executable)}: 1 for code segments, 0 for data.
    \item \textbf{Bit 1 (R/W)}: Readable (code) or writable (data) flag.
\end{itemize}

\subsection{The Far Jump and Segment Register Flush}

After loading the new GDT with \texttt{lgdt}, the CPU's \textit{segment descriptor
caches} (hidden registers that cache the properties of the current segments) still
hold the old BIOS values. Simply updating the data segment registers (\texttt{DS},
\texttt{ES}, etc.) with \texttt{mov} instructions is sufficient for those --- but
the \texttt{CS} (Code Segment) register cannot be updated with a \texttt{mov}.

The only way to load a new value into \texttt{CS} is via a \textbf{Far Jump}
(\texttt{jmp seg:offset}), which atomically updates \texttt{CS} and flushes the CPU's
instruction pipeline, forcing it to re-fetch instructions under the new segmentation:

\begin{lstlisting}[language={[x86masm]Assembler}, title=boot/gdt\_flush.asm]
    lgdt [eax]       ; Load the GDT pointer
    mov ax, 0x10     ; Selector 0x10 = 3rd descriptor = Kernel Data
    mov ds, ax
    mov es, ax
    jmp 0x08:.flush  ; Far jump: CS = selector 0x08 = Kernel Code
.flush:
    ret
\end{lstlisting}

\section{The Interrupt Descriptor Table (IDT)}

\subsection{The Interrupt Mechanism}

An \textbf{interrupt} is a hardware or software signal that causes the CPU to
temporarily suspend its current execution stream and transfer control to a specific
handler function. Interrupts are the fundamental mechanism through which hardware
communicates asynchronously with the software running on the CPU.

The x86 architecture defines 256 possible interrupt vectors (numbered 0--255). The
\textbf{IDT} is an array of 256 \textbf{gate descriptors}, each pointing to the
corresponding Interrupt Service Routine (ISR). The CPU's \texttt{IDTR} register
holds the base address and size of this table.

\subsection{Interrupt Vector Classification}

The 256 vectors are partitioned as follows:
\begin{itemize}
    \item \textbf{0--31}: Intel-reserved CPU Exceptions (e.g., \#DE Division by Zero,
    \#PF Page Fault, \#GP General Protection Fault). These cannot be remapped.
    \item \textbf{32--255}: Available for OS and hardware use. We map hardware IRQs
    here after remapping the PIC.
\end{itemize}

\subsection{The Gate Descriptor}

An IDT gate descriptor is 8 bytes and contains the full \texttt{selector:offset}
address of the ISR, plus a type/attribute byte. The \textbf{DPL field} in the
attribute byte is critical: it defines the minimum privilege level required to
trigger this interrupt via a software \texttt{int} instruction. Setting DPL to 3
(Ring 3) on interrupt 128 (\texttt{0x80}) allows User Mode programs to legally
invoke the kernel's system call handler.

\section{The Programmable Interrupt Controller (8259 PIC)}

\subsection{IRQ Routing}

Physical hardware devices (keyboard, mouse, disk controller, timer) communicate with
the CPU not by writing to memory, but by asserting an electrical signal on a dedicated
wire called an \textbf{Interrupt Request (IRQ) line}. The \textbf{8259A Programmable
Interrupt Controller} arbitrates these 15 hardware IRQs and delivers them to the
CPU as interrupt vectors.

The original IBM PC BIOS maps IRQ0--IRQ7 to CPU vectors 8--15. This collides directly
with Intel's reserved exception vectors (e.g., vector 8 is a Double Fault exception).
The OS \textit{must} remap the PIC to a safe range during initialization.

\subsection{PIC Initialization Sequence (ICW)}

The 8259 is programmed by sending a sequence of \textbf{Initialization Command Words
(ICW)} to its I/O ports. The master PIC lives at ports \texttt{0x20}/\texttt{0x21};
the slave at \texttt{0xA0}/\texttt{0xA1}.

\begin{enumerate}
    \item \textbf{ICW1} (\texttt{0x11} to command port): Start initialization, expect ICW4.
    \item \textbf{ICW2} (to data port): The new vector offset. We use 32 for master (IRQ0=32),
    40 for slave (IRQ8=40).
    \item \textbf{ICW3}: Describes the master/slave wiring topology.
    \item \textbf{ICW4} (\texttt{0x01}): Sets 8086/88 mode.
\end{enumerate}

After remapping, the PIC must send an \textbf{End of Interrupt (EOI)} signal
(\texttt{0x20} to port \texttt{0x20}) at the end of every IRQ handler, otherwise the
PIC will not deliver any further interrupts of equal or lower priority.

\section{CPU Exceptions and Kernel Panics}

The 32 CPU exceptions are the processor's way of signaling that something has gone
catastrophically wrong. The most important ones for OS development are:

\begin{itemize}
    \item \textbf{\#PF (14) --- Page Fault}: Triggered when a program accesses a virtual
    address that has no corresponding physical mapping in the Page Tables, or when it
    attempts to write to a read-only page. The faulting address is stored in the
    \texttt{CR2} register, which is invaluable for debugging.
    \item \textbf{\#GP (13) --- General Protection Fault}: Triggered by any attempt to
    violate x86 protection rules --- for example, a Ring 3 program trying to execute a
    privileged instruction like \texttt{cli} (disable interrupts) or accessing a
    segment it does not have permission for.
    \item \textbf{\#DF (8) --- Double Fault}: Triggered when the CPU encounters a fault
    \textit{while already processing a fault}. A Double Fault almost always indicates a
    corrupted kernel stack or a broken IDT.
\end{itemize}

For exceptions that push an \textbf{error code} onto the stack (e.g., \#PF, \#GP),
the ISR must pop this code before executing \texttt{iret}, otherwise the stack is
misaligned and the CPU will immediately triple-fault and reset.

\newpage

\chapter{Memory Management}

\section{Virtual Memory and Paging}

\subsection{The Isolation Problem}

In a flat physical memory model, every process can read and write any address in RAM.
A bug in one program can silently corrupt the kernel or another process's data. This
is fundamentally unsafe for a multi-process OS.

\textbf{Virtual Memory} solves this by introducing an abstraction layer between the
addresses a program uses (virtual addresses) and the actual physical RAM locations.
Every process sees a private, contiguous address space (e.g., always starting at
\texttt{0x00000000} and spanning 4GB), even though the underlying physical RAM is
fragmented and shared.

\subsection{The Two-Level Page Table Hierarchy}

x86 implements virtual memory through \textbf{Paging}. The 32-bit virtual address
space (4GB) is divided into \textbf{pages} of 4096 bytes each ($2^{12}$ bytes,
hence the bottom 12 bits of an address are the \textit{page offset}).

A naive approach would be a single 4GB/4KB = 1,048,576-entry lookup table --- but
storing 1 million 4-byte entries requires 4MB of RAM \textit{per process}, just for
the table itself. x86 solves this with a two-level hierarchy:

\begin{enumerate}
    \item \textbf{Page Directory (PD)}: 1024 entries, each pointing to a Page Table.
    The top 10 bits of the virtual address (bits 22--31) index into the PD.
    \item \textbf{Page Table (PT)}: 1024 entries, each pointing to a physical 4KB page.
    The next 10 bits (bits 12--21) index into the PT.
    \item \textbf{Page Offset}: The bottom 12 bits (bits 0--11) are added to the
    physical page base address to produce the final physical address.
\end{enumerate}

Address translation formula:
\[
\text{Physical Address} = \text{PT}[\text{PD}[\text{VAddr}[31{:}22]][VAddr[21{:}12]]] + VAddr[11{:}0]
\]

\subsection{Page Table Entry Format}

Each 4-byte Page Table Entry encodes both the physical page address and permission flags.
Since pages are always 4KB-aligned, the bottom 12 bits of the address are always zero,
so the hardware reuses them as flags:

\begin{itemize}
    \item \textbf{Bit 0 (Present)}: 1 = this page is mapped; 0 = accessing this page
    triggers a Page Fault (used to implement demand paging and swap).
    \item \textbf{Bit 1 (R/W)}: 0 = read-only; 1 = read-write.
    \item \textbf{Bit 2 (U/S)}: 0 = supervisor only (Ring 0); 1 = accessible by user
    space (Ring 3).
    \item \textbf{Bits 3--11}: Accessed, Dirty, Global, and other hardware flags.
    \item \textbf{Bits 12--31}: The upper 20 bits of the physical page frame address.
\end{itemize}

\subsection{Enabling Paging}

The Page Directory's physical address is loaded into the \textbf{CR3} control register.
Paging is then activated by setting \textbf{bit 31 of CR0} (the ``PG'' bit). The moment
this bit is set, \textit{every subsequent memory access} by the CPU is transparently
translated through the page tables. Identity mapping the first 4MB (virtual address =
physical address for the first 4MB) ensures the kernel does not crash immediately when
paging is enabled.

\section{Dynamic Memory: The Heap Allocator}

\subsection{Stack vs. Heap Allocation}

Local variables in C++ are allocated on the \textbf{stack} --- a fixed-size region of
memory (typically a few kilobytes to a few megabytes) that grows and shrinks with
function calls. Stack allocation is extremely fast (a single decrement of the \texttt{ESP}
register) but limited in size and lifetime (the memory is reclaimed when the function returns).

The \textbf{heap} is a pool of memory reserved for \textit{dynamic} allocation. Memory
on the heap persists until explicitly freed, regardless of which function allocated it.
This is necessary for data structures whose size is not known at compile time.

\subsection{The Bump Allocator}

The simplest possible heap allocator is the \textbf{bump allocator} (also called a
linear allocator or arena allocator). It maintains a single pointer to the ``watermark''
of the heap. Allocation simply advances this pointer by the requested size:

\begin{lstlisting}[language=C++, title=kernel/memory.cpp]
uint32_t heap_ptr = 0x200000; // Start of heap at 2MB mark

void* malloc(size_t size) {
    uint32_t addr = heap_ptr;
    heap_ptr += size;          // "Bump" the pointer forward
    return (void*)addr;
}
\end{lstlisting}

\textbf{Time complexity}: O(1) for allocation.\\
\textbf{Limitation}: \texttt{free()} is a no-op. Memory cannot be reclaimed.

Production kernels use more sophisticated allocators:
\begin{itemize}
    \item \textbf{Linked-list allocator}: Maintains a free-list of released blocks,
    merging adjacent free blocks (coalescing) to prevent fragmentation.
    \item \textbf{Buddy allocator}: Splits memory into power-of-two blocks; merging
    is efficient because a block's ``buddy'' is at a predictable offset.
    \item \textbf{Slab allocator} (Linux): Pre-allocates pools of fixed-size objects
    to minimize fragmentation for frequent kernel allocations like task structures.
\end{itemize}

\newpage

\chapter{Kernel Services}

\section{The Interactive Shell}

A kernel without user interaction is a black box. The \textbf{command-line shell}
is the most fundamental interactive interface between the user and the OS. It requires
three building blocks: keyboard input (via the IRQ handler), a string buffer to
accumulate keystrokes, and a command dispatcher that maps strings to functions.

The keyboard IRQ (IRQ 1, vector 33) fires on every key press and release. The handler
reads a \textbf{scancode} from I/O port \texttt{0x60} --- a hardware-specific byte
indicating which physical key changed state. A \textbf{scancode-to-ASCII lookup table}
maps these scancodes to printable characters.

\section{Memory Segmentation and Ring Isolation: The TSS}

\subsection{Privilege Rings}

x86 defines four privilege levels, called \textbf{Rings}, numbered 0--3. Ring 0 is the
most privileged (the kernel); Ring 3 is the least privileged (user applications).
The current privilege level (CPL) is stored in bits 0--1 of the CS segment register.

Hardware enforces this separation: an instruction like \texttt{cli} (disable interrupts)
executed in Ring 3 immediately triggers a \#GP exception. This prevents a buggy or
malicious application from disabling the timer interrupt and freezing the OS.

\subsection{The Task State Segment (TSS)}

When a hardware interrupt fires in Ring 3, the CPU must switch to the kernel stack
(Ring 0 stack) to handle it securely. Otherwise, a user program could manipulate the
interrupt handler by crafting a fake stack frame. The CPU finds the Ring 0 stack
address by reading it from the \textbf{Task State Segment (TSS)}.

The TSS is a special data structure (minimum 104 bytes) registered in the GDT under a
\textbf{system descriptor} (S-bit = 0, type = \texttt{0x9}). The TSS descriptor selector
is loaded into the \texttt{TR} (Task Register) via the \texttt{ltr} instruction. When a
privilege transition occurs, the CPU automatically reads \texttt{ESP0} and \texttt{SS0}
from the TSS to set up the kernel stack.

\section{User Mode (Ring 3)}

\subsection{Dropping Privileges}

The CPU always starts in Ring 0. To run untrusted user code, we must artificially
transfer execution to Ring 3. This is done by constructing a \textbf{fake interrupt
return frame} on the stack and executing \texttt{iret}.

When \texttt{iret} executes, it pops \texttt{EIP}, \texttt{CS}, \texttt{EFLAGS}, \texttt{ESP},
and \texttt{SS} from the stack. If the \texttt{CS} selector's DPL is 3, the CPU
transitions to Ring 3 and loads the user-mode stack from the popped \texttt{SS:ESP}.

\begin{lstlisting}[language=C++, title=boot/usermode.asm -- Fake IRET frame for Ring 3]
jump_usermode:
    mov ax, 0x23    ; User data segment selector (GDT index 4, RPL=3)
    mov ds, ax
    push 0x23       ; SS
    push user_stack ; ESP
    pushf           ; EFLAGS (current flags, including IF=1)
    push 0x1B       ; CS = user code segment (GDT index 3, RPL=3)
    push user_main  ; EIP (entry point of user program)
    iret            ; CPU drops to Ring 3!
\end{lstlisting}

\section{System Calls: The Kernel--Userspace Bridge}

\subsection{The Fundamental Isolation Problem}

Ring 3 code is physically incapable of accessing kernel services directly. A user
program cannot write to VGA memory at \texttt{0xB8000} (mapped read-only for Ring 3),
call kernel functions (different address space), or execute privileged instructions.

Without a mechanism to cross this boundary, user programs are entirely useless.

\subsection{The \texttt{int 0x80} ABI}

The universal solution is the \textbf{system call}. A user program places arguments
in CPU registers and triggers a predefined software interrupt. The CPU performs a
\textit{controlled} privilege escalation: it validates the IDT gate's DPL (must be
$\geq$ Ring 3 to be callable from Ring 3), switches to the kernel stack via the TSS,
and jumps to the kernel's interrupt handler at Ring 0.

The calling convention for Abyss system calls uses CPU registers to pass arguments
(since the user and kernel stacks are separate and cannot easily share pointers
across privilege boundaries):
\begin{itemize}
    \item \textbf{EAX}: System call number (1 = print string).
    \item \textbf{EBX}: First argument (e.g., pointer to string in user memory).
\end{itemize}

The kernel handler inspects \texttt{EAX} to dispatch the correct service, performs
the operation on behalf of the user program, and returns via \texttt{iret}.

\section{The System Timer (PIT) and Preemptive Multitasking}

\subsection{The Programmable Interval Timer}

The \textbf{8253/8254 Programmable Interval Timer (PIT)} is a hardware chip that
generates periodic interrupts at a configurable frequency. Its internal oscillator
runs at exactly \textbf{1,193,182 Hz}. To achieve a tick rate of $f$ Hz, we program
a divisor $d$:
\[
d = \left\lfloor \frac{1{,}193{,}182}{f} \right\rfloor
\]
For $f = 100$ Hz (100 ticks/second): $d = 11{,}931$.

The divisor is written to the PIT's data port (\texttt{0x40}) in two 8-bit writes
(low byte first, then high byte). The PIT is wired to IRQ 0, so it fires vector 32 100 times per second.

\subsection{Multitasking Theory: The Illusion of Concurrency}

A single-core CPU executes exactly one instruction at a time. \textbf{Multitasking}
creates the \textit{illusion} of parallel execution by rapidly switching between tasks,
each of which perceives it has sole ownership of the CPU.

The OS maintains a \textbf{Task Control Block (TCB)} for each task, storing its complete
CPU state (registers, stack pointer, instruction pointer) when it is descheduled.

\textbf{Cooperative multitasking} requires tasks to voluntarily yield the CPU
(\texttt{yield()} call). A single hung task freezes the entire system.

\textbf{Preemptive multitasking} uses the PIT interrupt to forcibly switch tasks
100 times per second, regardless of task behavior. The context switch happens inside
the IRQ 0 handler:

\begin{enumerate}
    \item IRQ 0 fires. CPU pushes \texttt{EFLAGS}, \texttt{CS}, \texttt{EIP} onto the
    \textit{current task's stack} (this is the hardware-assisted part of the state save).
    \item The \texttt{isr32} handler pushes all general-purpose registers (\texttt{pusha}).
    \item The current \texttt{ESP} is saved into the current TCB.
    \item \texttt{ESP} is overwritten with the next task's saved \texttt{ESP} from its TCB.
    \item Registers are restored from the new stack (\texttt{popa}).
    \item \texttt{iret} executes, which pops the new task's \texttt{EIP}/\texttt{CS}/\texttt{EFLAGS}
    and resumes it \textit{exactly where it was paused}.
\end{enumerate}

The power of this mechanism is that the context switch is entirely \textit{transparent}
to the tasks --- they are never aware they were paused.

\newpage

\chapter{Virtual File System}

\section{The Abstraction Problem}

A fully functional OS must manage persistent data. However, there are dozens of
incompatible storage formats: FAT32, ext4, NTFS, APFS, ZFS. If every application
had to know the details of the underlying file system, porting an application from
one OS to another would require rewriting all its storage code.

The \textbf{Virtual File System (VFS)} is an abstraction layer that presents a
unified API to applications, regardless of the underlying storage driver. The same
\texttt{read()} call works whether the file lives on a FAT32 USB drive, an ext4
hard drive, or in RAM.

\section{The inode Model}

The VFS models every storage object as a \textbf{node} (analogous to a Unix
\texttt{inode} --- index node). Each node contains metadata about a file (name, size,
type) plus a \textbf{function pointer} to the driver-specific read implementation:

\begin{lstlisting}[language=C++, title=kernel/vfs.h]
struct fs_node {
    char     name[128];
    uint32_t length;
    uint32_t flags;   // VFS_FILE or VFS_DIRECTORY
    uint32_t (*read)(fs_node*, uint32_t offset, uint32_t size, char* buf);
};
\end{lstlisting}

The function pointer is the key to polymorphism in C. A FAT32 node's \texttt{read}
pointer points to \texttt{fat32\_read()}; a TarFS node's points to \texttt{tarfs\_read()}.
The kernel's generic \texttt{vfs\_read()} simply calls \texttt{node->read()} without
knowing which driver is underneath.

\section{The Initial Ramdisk (initrd)}

Before a real disk driver exists, we need a way to provide initial files to the kernel.
The solution is the \textbf{Initial Ramdisk}: a filesystem image that the bootloader
loads entirely into RAM before handing control to the kernel.

GRUB supports loading \textit{modules} (binary blobs) into RAM and passing their
physical addresses to the kernel via the Multiboot Information Structure
(\texttt{mods\_addr} field). We package files into a \textbf{ustar TAR archive}
and instruct GRUB to load it as a module.

\subsection{The ustar TAR Format}

The TAR format is elegantly simple. A TAR archive is a flat sequence of 512-byte
\textit{blocks}. Each file is preceded by a \textbf{512-byte header block} containing:
\begin{itemize}
    \item \textbf{filename[100]}: The file name (null-terminated).
    \item \textbf{filesize[12]}: The file size in \textit{octal ASCII} (e.g.,
    \texttt{"00000000017"} for 15 bytes). Storing the size in octal ASCII is a
    historical quirk: early UNIX ran on PDP machines that natively used octal.
    \item \textbf{magic[6]}: \texttt{"ustar"} for compliance validation.
\end{itemize}
Following the header, the raw file data occupies the next $\lceil size/512 \rceil$
blocks, zero-padded to a 512-byte boundary.

The TarFS driver walks through these blocks in memory, builds an \texttt{fs\_node}
for each file, and wires the function pointer to \texttt{tarfs\_read}, which simply
copies bytes from the in-RAM location.

\newpage

\chapter{Graphics \& Graphical User Interface}

\section{VESA, VBE, and the Linear Framebuffer}

\subsection{Beyond Text Mode}

VGA Text Mode's character-grid paradigm is fundamentally incapable of rendering
arbitrary graphics. A pixel-level graphical display requires a different model:
the \textbf{Linear Framebuffer (LFB)}.

\textbf{VESA BIOS Extensions (VBE)} is a standard API (originally implemented in
BIOS firmware, later superseded by UEFI GOP) that allows software to query the
GPU for available video modes and switch into one. In Multiboot mode, we request a
video mode via the Multiboot header flags, and GRUB configures VBE on our behalf
before booting the kernel.

\subsection{The Linear Framebuffer Model}

In a linear framebuffer, the GPU exposes its video RAM as a contiguous array of
pixels in the CPU's physical address space. In \texttt{1024$\times$768} with 32-bit
color, the framebuffer is:
\[
1024 \times 768 \times 4 \text{ bytes} = 3{,}145{,}728 \text{ bytes} \approx 3 \text{ MB}
\]

The address of pixel $(x, y)$ is:
\[
\text{addr} = \text{fb\_base} + y \times \text{pitch} + x \times 4
\]

where \textbf{pitch} (or stride) is the number of bytes per row. Pitch may be larger
than $\text{width} \times 4$ if the GPU pads rows for alignment.

Each 32-bit pixel encodes color as \texttt{0x00RRGGBB}: one byte each for red, green,
and blue components, with the top byte unused (or used for alpha transparency in some
modes).

\section{Double Buffering}

\subsection{The Tearing Problem}

If we draw directly to the framebuffer (which the display hardware is simultaneously
reading to refresh the screen), we risk a visual artifact called \textbf{screen tearing}:
the monitor displays the top half of the old frame and the bottom half of the new frame
simultaneously, producing a horizontal seam.

More critically, drawing a complex scene takes time. If the cursor is drawn before the
background is cleared, intermediate states become visible, causing flicker.

\subsection{The Double-Buffer Solution}

\textbf{Double buffering} introduces a second, identical-sized memory region called the
\textbf{back buffer}. All rendering operations target the back buffer exclusively. When
the entire scene is finished, the contents of the back buffer are copied to the
\textbf{front buffer} (the actual framebuffer) in a single, fast operation:

\begin{lstlisting}[language=C++, title=kernel/vesa.cpp]
void vesa_swap_buffers() {
    // Copy 3MB of backbuffer to framebuffer in one blast
    uint32_t dwords = (fb_width * fb_height);
    asm volatile (
        "rep movsl"
        : "+S"(backbuffer), "+D"(framebuffer), "+c"(dwords)
        : : "memory"
    );
}
\end{lstlisting}

The \texttt{rep movsl} assembly instruction copies \texttt{ECX} double-words (4-byte
chunks) from \texttt{ESI} to \texttt{EDI} at maximum CPU throughput. Copying 3MB this
way is orders of magnitude faster than a C++ loop, as the CPU can use internal
write-combining buffers.

\section{Rasterization Algorithms}

\subsection{Linear Interpolation for Gradients}

A gradient smoothly transitions between two colors. For a vertical gradient between
color $c_1$ (top) and $c_2$ (bottom), row $i$ out of $h$ total rows receives the
interpolated color computed independently for each channel:
\[
R_i = R_1 + \frac{(R_2 - R_1) \cdot i}{h}, \quad G_i = G_1 + \frac{(G_2 - G_1) \cdot i}{h}, \quad B_i = B_1 + \frac{(B_2 - B_1) \cdot i}{h}
\]

This is \textbf{linear interpolation (lerp)}, the simplest form of color blending.
The division is integer division (no floating point required), performed using the
full-precision intermediate value before truncation.

\subsection{Bresenham's Line Algorithm}

Drawing a line between two arbitrary points requires determining which pixels best
approximate the mathematical line. Floating-point math (slope, $y = mx + b$) is too
slow for per-pixel computation in a kernel.

\textbf{Bresenham's algorithm} (1962) uses only integer arithmetic. It maintains an
\textbf{error accumulator} $e$ that tracks the deviation of the rasterized line from
the ideal. At each step, the error is compared to the slope to decide whether to
advance both $x$ and $y$, or only $x$:

\begin{lstlisting}[language=C++]
int dx = abs(x1 - x0), dy = abs(y1 - y0);
int err = dx + dy;
while (true) {
    put_pixel(x0, y0, color);
    if (x0 == x1 && y0 == y1) break;
    int e2 = 2 * err;
    if (e2 >= dy) { err += dy; x0 += sx; }
    if (e2 <= dx) { err += dx; y0 += sy; }
}
\end{lstlisting}

\subsection{Midpoint Circle Algorithm}

A circle of radius $r$ satisfies $x^2 + y^2 = r^2$. The \textbf{midpoint circle
algorithm} exploits the circle's 8-fold symmetry: computing one octant (45 degrees)
yields all 8 symmetric points simultaneously. It also avoids square roots by
maintaining an integer \textbf{decision variable} $d$ initialized to $3 - 2r$. At
each step, $d$ determines whether to move to the next row, and the decision variable
is updated with only additions and subtractions.

\subsection{Bitmap Font Rendering}

Switching to VESA graphics mode disables the VGA hardware character generator.
Text must be rendered manually using a \textbf{bitmap font}: a lookup table where
each character is represented as an $8 \times 8$ grid of binary pixels.

For character ASCII code $c$ and glyph row $r$, the 8-bit value \texttt{font[c][r]}
encodes which of the 8 pixels in that row should be drawn. Pixel $p$ (0--7) is
extracted with a bitwise shift and mask:
\[
\text{draw pixel}(p, r) \iff \left(\texttt{font}[c][r] \gg p\right) \mathbin{\&} 1 = 1
\]

By replacing \texttt{put\_pixel} with \texttt{draw\_rect(x, y, scale, scale)}, the
kernel can scale the $8 \times 8$ bitmap to any integer multiple (e.g., $\times 2$
for $16 \times 16$ text) without a separate font atlas for each size.

\section{The PS/2 Mouse: Hardware Protocol}

\subsection{The PS/2 Controller (Intel 8042)}

The PS/2 keyboard and mouse share a legacy controller chip (the Intel 8042 or
compatible) with two channels:
\begin{itemize}
    \item \textbf{Primary port} (I/O port \texttt{0x60}/\texttt{0x64}): Keyboard.
    \item \textbf{Auxiliary port}: Mouse. Must be explicitly enabled by setting bit 1
    of the controller's configuration byte.
\end{itemize}

\subsection{The 3-Byte Packet Format}

The PS/2 mouse transmits movement data in 3-byte packets, one byte per IRQ 12 firing.
A state machine in the handler assembles the packet:

\begin{center}
\begin{tabular}{|c|c|c|}
\hline
\textbf{Byte} & \textbf{Bits} & \textbf{Meaning} \\
\hline
0 & [0] & Left button pressed \\
0 & [1] & Right button pressed \\
0 & [3] & Synchronization bit (always 1 in valid packets) \\
0 & [4] & X sign bit (extends X delta to 9-bit signed) \\
0 & [5] & Y sign bit (extends Y delta to 9-bit signed) \\
1 & [7:0] & X movement delta (8-bit, sign-extended using byte 0 bit 4) \\
2 & [7:0] & Y movement delta (inverted: positive = down in hardware coords) \\
\hline
\end{tabular}
\end{center}

The Y-axis inversion is critical: the PS/2 protocol uses a coordinate system where
positive Y is upward, while screen coordinates have positive Y downward. The driver
must subtract (not add) the Y delta from the cursor's screen Y coordinate.

\section{Window Management and the Compositor}

\subsection{The Compositor Pattern}

A \textbf{compositor} is the subsystem responsible for assembling the final scene from
multiple independently-rendered layers and presenting it on screen. In Abyss OS, the
\texttt{ui\_render()} function acts as a simple compositor:

\begin{enumerate}
    \item \textbf{Clear}: Render the desktop background gradient to the back buffer.
    \item \textbf{Paint Windows}: Iterate through the window array (back-to-front,
    respecting Z-order) and render each window's chrome (title bar, borders, body)
    plus its application-specific content.
    \item \textbf{Draw Cursor}: Render the mouse cursor on top of all windows.
    \item \textbf{Swap}: Atomically copy the back buffer to the front buffer.
\end{enumerate}

\subsection{Z-Order and Window Focus}

The \textbf{Z-order} determines which window appears on top when windows overlap.
In Abyss, the \texttt{windows[]} array is ordered back-to-front: the last element
is always rendered on top. When the user clicks a window, it is \textit{promoted}:
removed from its current position, and appended to the end of the array using a
shift operation, ensuring it is rendered over all other windows.

\subsection{Click Detection and Drag Mechanics}

Drag detection requires distinguishing between a fresh click and a held button.
The key insight is: a drag should only \textit{begin} on the exact frame where
\texttt{mouse\_left\_pressed} transitions from \texttt{false} to \texttt{true}.
This is tracked by comparing the current state against \texttt{prev\_mouse\_left\_pressed},
a boolean updated at the end of each frame.

A window is \textit{draggable} only if the click started within its title bar. During
a drag, the window's position is updated to:
\[
\text{window.x} = \text{mouse.x} - \text{drag\_offset\_x}
\]
where \texttt{drag\_offset\_x} was recorded as the distance between the cursor and the
window's top-left corner \textit{at the moment the drag began}, ensuring the window
does not snap its origin to the cursor.

\newpage

\chapter{Storage \& FAT32 File System}

\section{ATA and Disk I/O Theory}

\subsection{The Physical Hard Drive}

A \textbf{Hard Disk Drive (HDD)} stores data by magnetically encoding bits onto
rotating metallic platters. Data is organized into concentric \textbf{tracks}; each
track is divided into \textbf{sectors} of 512 bytes. A \textbf{read/write head}
floats nanometers above the platter surface on an actuator arm, reading or writing
bits as the platter spins beneath it.

\textbf{Seek time} (time to position the head over the correct track) and
\textbf{rotational latency} (time to wait for the target sector to rotate under the
head) are the two dominant sources of latency for HDD reads, typically totaling
5--15 ms --- approximately 10,000,000 CPU cycles at 1 GHz.

\subsection{Logical Block Addressing (LBA)}

Early BIOSes addressed disk sectors using the physical \textbf{CHS (Cylinder-Head-Sector)}
coordinate system. This required the OS to know the physical geometry of each drive ---
an impractical dependency as drives evolved.

\textbf{Logical Block Addressing (LBA)} abstracts the disk into a flat, zero-indexed
array of 512-byte blocks. LBA28 uses a 28-bit address, supporting:
\[
2^{28} \times 512 = 137{,}438{,}953{,}472 \text{ bytes} = 128 \text{ GB}
\]

\subsection{ATA PIO Mode}

\textbf{ATA (Advanced Technology Attachment)} defines the electrical and logical
interface between the motherboard and storage devices. The legacy \textbf{PIO
(Programmed I/O)} mode transfers data by having the CPU manually execute I/O
instructions for every word of data. While slower than DMA (Direct Memory Access),
PIO is the simplest baseline and requires no additional hardware support.

The ATA controller lives at I/O ports \texttt{0x1F0}--\texttt{0x1F7} (primary bus):

\begin{center}
\begin{tabular}{|c|c|}
\hline
\textbf{Port} & \textbf{Purpose} \\
\hline
\texttt{0x1F0} & Data register (16-bit, read/write 1 sector = 256 words) \\
\texttt{0x1F1} & Error register (read) / Features (write) \\
\texttt{0x1F2} & Sector count \\
\texttt{0x1F3} & LBA bits 0--7 \\
\texttt{0x1F4} & LBA bits 8--15 \\
\texttt{0x1F5} & LBA bits 16--23 \\
\texttt{0x1F6} & Drive select + LBA bits 24--27 \\
\texttt{0x1F7} & Status (read) / Command (write) \\
\hline
\end{tabular}
\end{center}

A read sequence:
\begin{enumerate}
    \item Write LBA address across ports \texttt{0x1F3}--\texttt{0x1F6}.
    \item Write command \texttt{0x20} (READ SECTORS) to \texttt{0x1F7}.
    \item \textbf{Poll} status port \texttt{0x1F7}: wait until \textbf{BSY} (bit 7)
    clears and \textbf{DRQ} (bit 3, Data Request) sets.
    \item Execute \texttt{rep insw} to read 256 16-bit words from port \texttt{0x1F0}
    into the buffer.
\end{enumerate}

\section{The FAT32 File System}

\subsection{Volume Layout}

A FAT32-formatted disk is divided into four regions:

\begin{center}
\begin{tabular}{|c|c|c|}
\hline
\textbf{Region} & \textbf{Size} & \textbf{Purpose} \\
\hline
VBR (Sector 0) & 1 sector & Volume metadata (BPB) \\
Reserved Sectors & Variable & GRUB, FSInfo, backup boot sector \\
FAT Table(s) & Variable & Cluster allocation linked list \\
Data Region & Remainder & Actual file and directory content \\
\hline
\end{tabular}
\end{center}

\subsection{The Volume Boot Record and BIOS Parameter Block}

The \textbf{Volume Boot Record (VBR)}, also called the Partition Boot Record, occupies
the first 512 bytes of the partition. It contains the \textbf{BIOS Parameter Block (BPB)},
a structured header describing the file system geometry.

Key BPB fields relevant to our implementation:
\begin{itemize}
    \item \textbf{bytes\_per\_sector} (offset 11): Almost always 512.
    \item \textbf{sectors\_per\_cluster} (offset 13): Power of 2 (1, 2, 4, ..., 128).
    Larger values reduce FAT table size but increase internal fragmentation.
    \item \textbf{reserved\_sectors} (offset 14): Number of sectors before the FAT
    (where the VBR lives). Typically 32 for FAT32.
    \item \textbf{fat\_count} (offset 16): Number of FAT copies (usually 2 for redundancy).
    \item \textbf{sectors\_per\_fat} (offset 36): Size of one FAT table in sectors.
    \item \textbf{root\_cluster} (offset 44): Cluster number of the Root Directory.
\end{itemize}

From these values, the LBA of any region can be computed:
\[
\text{fat\_lba} = \text{reserved\_sectors}
\]
\[
\text{data\_lba} = \text{fat\_lba} + (\text{fat\_count} \times \text{sectors\_per\_fat})
\]
\[
\text{cluster\_lba}(n) = \text{data\_lba} + (n - 2) \times \text{sectors\_per\_cluster}
\]

The $-2$ offset is because FAT32 cluster numbering starts at 2 (clusters 0 and 1
are reserved for special markers).

\subsection{The File Allocation Table}

The \textbf{File Allocation Table} is an array of 32-bit unsigned integers (28 bits
used, top 4 reserved). Each entry corresponds to a cluster in the Data Region:
\begin{itemize}
    \item \texttt{0x00000000}: Cluster is free.
    \item \texttt{0x00000002}--\texttt{0x0FFFFFEF}: Next cluster in the file's chain.
    \item \texttt{0x0FFFFFF8}--\texttt{0x0FFFFFFF}: End-of-File (EOF) marker.
    \item \texttt{0x0FFFFFF7}: Bad cluster (hardware defect, skip).
\end{itemize}

Files are stored in a \textbf{singly-linked list of clusters}. To read a file starting
at cluster $C_0$, the driver follows the chain: $C_0 \to \text{FAT}[C_0] \to \text{FAT}[\text{FAT}[C_0]] \to \ldots \to \text{EOF}$.

To look up entry $C$ in the FAT, the driver must:
\begin{enumerate}
    \item Compute the FAT sector: $\text{fat\_sector} = \text{fat\_lba} + \lfloor C / 128 \rfloor$ (since $512/4 = 128$ entries per sector).
    \item Compute the byte offset within that sector: $\text{offset} = (C \bmod 128) \times 4$.
    \item Read the sector via ATA, extract the 4 bytes at the offset, mask with
    \texttt{0x0FFFFFFF}.
\end{enumerate}

\subsection{Directory Entries}

Directories in FAT32 are just files whose content is an array of \textbf{32-byte
directory entries}. Each entry describes one file or subdirectory:

\begin{center}
\begin{tabular}{|c|c|c|}
\hline
\textbf{Offset} & \textbf{Size} & \textbf{Field} \\
\hline
0   & 8 bytes & Short filename (space-padded, uppercase) \\
8   & 3 bytes & Extension (space-padded, uppercase) \\
11  & 1 byte  & Attributes bitmask \\
20  & 2 bytes & First cluster high word \\
26  & 2 bytes & First cluster low word \\
28  & 4 bytes & File size in bytes \\
\hline
\end{tabular}
\end{center}

The \textbf{attributes byte} is a bitmask:
\begin{itemize}
    \item \texttt{0x01}: Read-only.
    \item \texttt{0x02}: Hidden.
    \item \texttt{0x04}: System file.
    \item \texttt{0x08}: Volume label (not a real file --- skip it).
    \item \texttt{0x10}: Directory.
    \item \texttt{0x0F}: Long File Name entry (skip; handled by LFN subsystem).
\end{itemize}

A special first-byte value of \texttt{0xE5} marks a \textit{deleted} entry
(the disk space is free but the content is not yet overwritten). A value of
\texttt{0x00} indicates that this entry and all subsequent entries in the directory
are empty --- the traversal can stop.

\subsection{Structure Alignment: \texttt{pragma pack}}

C++ compilers insert \textbf{padding bytes} between struct fields to align them on
natural boundaries (e.g., a 4-byte \texttt{uint32\_t} is placed at a 4-byte-aligned
offset). This optimization improves CPU memory access speed but means a C++ struct
may not match the physical byte layout on disk.

For example, without packing, a struct with a \texttt{char} followed by a
\texttt{uint32\_t} would have 3 invisible padding bytes between them, shifting all
subsequent fields by 3 bytes and producing garbage when cast from raw disk data.

\texttt{\#pragma pack(push, 1)} instructs the compiler to use 1-byte alignment
(no padding) for all subsequent structs. After the FAT32 structs, \texttt{\#pragma pack(pop)}
restores the default alignment. This ensures the C++ struct is a perfect byte-for-byte
mirror of the on-disk layout.

\newpage

\chapter{Interactive File Reading}

\section{Traversing the Full Root Directory}

\subsection{The Single-Sector Limitation}

The initial implementation of \texttt{fat32\_get\_root\_entries()} read only the first
sector of the root directory cluster. Since each 32-byte directory entry occupies one
slot in the 512-byte sector, the first sector contains exactly 16 entries. With Long
File Name (LFN) entries consuming additional slots (one LFN entry per 13 characters
of the long name), the effective limit on visible files was even lower.

The correct implementation must traverse the root directory's entire cluster chain,
just as \texttt{fat32\_read\_file()} traverses a regular file's chain:

\begin{lstlisting}[language=C++, title=kernel/fat32.cpp]
int fat32_get_root_entries(FAT32_DirectoryEntry* out_buffer, int max_entries) {
    uint32_t current_cluster = root_dir_cluster;
    int found_count = 0;
    while (current_cluster < 0x0FFFFFF8) {
        uint32_t lba = cluster_to_lba(current_cluster);
        for (int s = 0; s < sectors_per_cluster; s++) {
            uint8_t buf[512];
            ata_read_sector(lba + s, buf);
            FAT32_DirectoryEntry* entries = (FAT32_DirectoryEntry*)buf;
            for (int i = 0; i < 16; i++) {
                if (entries[i].name[0] == 0x00) return found_count;
                if ((uint8_t)entries[i].name[0] == 0xE5) continue;
                if (entries[i].attributes == 0x0F ||
                    (entries[i].attributes & 0x08)) continue;
                if (found_count < max_entries)
                    out_buffer[found_count++] = entries[i];
            }
        }
        current_cluster = fat32_get_next_cluster(current_cluster);
    }
    return found_count;
}
\end{lstlisting}

\section{Static Caching in the GUI}

Reading from the ATA controller is expensive: each sector read requires a polling loop
(busy-waiting on the BSY and DRQ status bits) that may stall the CPU for millions of
cycles. If \texttt{fat32\_get\_root\_entries()} were called on every frame of the GUI
render loop (typically 30--60 times per second), the system would spend the majority of
its time waiting on disk I/O, making window dragging jerky and unresponsive.

The solution is a \textbf{static cache}. In C++, a \texttt{static} local variable is
initialized only once and persists between function calls. Using a sentinel value
(\texttt{-1}) to indicate ``not yet loaded'', the file list is read from disk exactly
once and stored in a static buffer for all subsequent frames:

\begin{lstlisting}[language=C++]
static FAT32_DirectoryEntry files[16];
static int file_count = -1;
if (file_count == -1)
    file_count = fat32_get_root_entries(files, 16);
\end{lstlisting}

\section{Event-Driven Click Detection}

\subsection{The Hit-Test Problem}

Determining whether a mouse click landed on a specific UI element is called a
\textbf{hit test}. For rectangular elements (which all our file entries are), the
hit test is an axis-aligned bounding-box intersection:
\[
\text{hit} \iff x_{mouse} \in [x_{elem},\ x_{elem} + w_{elem}] \;\wedge\; y_{mouse} \in [y_{elem},\ y_{elem} + h_{elem}]
\]

Since each file entry is rendered at a dynamically computed Y position
(\texttt{file\_y} is incremented by 24 pixels per entry), the hit test must be
evaluated \textit{inside the rendering loop}, where \texttt{file\_y} has the correct
value for the current entry.

\subsection{Edge Detection for Mouse Clicks}

A naive implementation using \texttt{if (mouse\_left\_pressed)} would trigger
continuously for every frame the button is held. This would cause a single click on a
file to load it dozens of times and fire spuriously as the user drags the mouse across
files.

The correct approach is \textbf{rising-edge detection}: we only act on the first frame
where the button transitions from \textit{not pressed} to \textit{pressed}. This is
implemented by storing the previous frame's button state:
\[
\text{click} \iff \texttt{mouse\_left\_pressed} \wedge \neg\texttt{prev\_mouse\_left\_pressed}
\]

The previous state is updated at the end of every frame.

\section{Window Visibility and the Application Launch Model}

\subsection{Visibility as State}

A \textbf{visibility flag} on the \texttt{Window} struct allows windows to be
pre-allocated but hidden. This avoids the overhead of dynamic memory allocation for
window creation and simplifies the window management logic.

The Notepad window is pre-allocated with \texttt{visible = false}. When the user clicks
a file in the File Explorer, the OS:
\begin{enumerate}
    \item Reads the file's cluster chain from the FAT32 hard drive into a shared static buffer.
    \item Sets the Notepad window's \texttt{visible} flag to \texttt{true}.
    \item Sets \texttt{needs\_redraw = true} to force a compositor refresh.
\end{enumerate}

On the next frame, the compositor detects the visible Notepad window, renders its
white interior, and calls \texttt{draw\_string} with the shared buffer to rasterize
the file's text content.

\begin{lstlisting}[language=C++, title=kernel/ui.cpp -- Notepad rendering]
// Is this the Notepad window?
if (w->title[0] == 'N' && w->title[1] == 'o') {
    draw_rect(w->x+2, w->y+24, w->width-4, w->height-26, 0xFFFFFF);
    draw_string(notepad_text, w->x+10, w->y+30, 0x000000, 0xFFFFFF, 1);
}
\end{lstlisting}

\newpage


\chapter{Executable Programs \& The Loader API}

A defining characteristic of an operating system is its ability to load and execute standalone programs stored on secondary storage. In Abyss OS, this mechanism bridges the FAT32 file system, the graphical user interface, and the physical CPU execution pipeline.

\section{The Flat Binary Format}
Executable files typically rely on structured formats such as ELF (Executable and Linkable Format) or PE (Portable Executable), which contain complex headers, symbol tables, and relocation data. While these formats offer flexibility, they introduce significant parsing complexity.

For its initial implementation, Abyss OS utilizes the \textbf{Flat Binary} format. A flat binary contains raw machine code bytes starting exactly at offset zero, with no headers. The compiler and linker are instructed to assume the program will be loaded at a fixed virtual address (in our case, \texttt{0x400000}). Because our paging system maps the first 4GB using 4MB Big Pages, this address is already safely mapped and accessible in Ring 0.

To produce a flat binary, the cross-compiler is invoked with specific flags:
\begin{itemize}
    \item \texttt{-ffreestanding -nostdlib}: Prevents linking against standard C libraries.
    \item \texttt{-Wl,-Ttext,0x400000,-e,main}: Instructs the linker to resolve all internal memory references relative to \texttt{0x400000} and sets \texttt{main} as the entry point.
    \item \texttt{-Wl,--oformat,binary}: Strips all ELF metadata, resulting in pure x86 instructions.
\end{itemize}

\section{The Loader Mechanism}
The loader is responsible for transferring the program from the FAT32 filesystem into physical memory and jumping to it.

\begin{lstlisting}[language=C++, title=kernel/loader.cpp -- The Loader]
#define USER_LOAD_ADDRESS 0x400000

void load_and_run(uint8_t* binary, uint32_t size) {
    uint8_t* dest = (uint8_t*)USER_LOAD_ADDRESS;
    for (uint32_t i = 0; i < size; i++) {
        dest[i] = binary[i];
    }

    // Cast the load address to a function pointer
    typedef void (*program_t)(void (*)(const char*));
    program_t program = (program_t)USER_LOAD_ADDRESS;
    
    // Execute the program!
    program(sys_open_notepad);
}
\end{lstlisting}

\section{The Kernel API: Bypassing Syscalls for Ring 0 Execution}
Modern operating systems rely on software interrupts (e.g., \texttt{int 0x80}) to transition the CPU from User Mode (Ring 3) to Kernel Mode (Ring 0) when a program requests a service. However, because our initial programs run in Kernel Mode, using a software interrupt causes stack frame misalignment due to the absence of privilege-level switching (the CPU omits pushing \texttt{ESP} and \texttt{SS}).

To solve this, Abyss OS provides an API via direct C++ function pointers. The kernel hands the executing program a pointer to a kernel-level function, which the program can call identically to a local function.

\begin{lstlisting}[language=C, title=programs/hello.c -- The User Program]
// Define a function pointer type matching the kernel API
typedef void (*api_open_notepad)(const char*);

void main(api_open_notepad open_notepad) {
    // Call the kernel API directly
    open_notepad("SUCCESS!\n\nThis text was written by HELLO.BIN.\n"
                 "I am a standalone C program loaded\n"
                 "off the FAT32 hard drive!");
}
\end{lstlisting}

When the program executes \texttt{open\_notepad}, it triggers \texttt{sys\_open\_notepad} in the kernel, copying the requested string into the GUI's shared \texttt{notepad\_text} buffer and forcing the Notepad window to render. This effectively demonstrates code execution and cross-boundary communication.

\newpage
\appendix

\chapter{Debugging Journal: Errors Encountered and Solutions}

This appendix documents every major bug encountered during the development of Abyss OS, the symptoms that revealed them, and the exact fixes applied.

This appendix documents every major bug encountered during the development of Abyss OS, the symptoms that revealed them, and the exact fixes applied. It serves as a practical reference for common bare-metal programming mistakes.

% ─────────────────────────────────────────────

\subsection*{Phase 27: Executable Freezes \& the ELF Header Trap}
\textbf{Symptom:} Double-clicking \texttt{HELLO.BIN} in the File Explorer caused the OS to instantly freeze (the mouse cursor stopped moving). \\
\textbf{Diagnosis:} Although the file was named \texttt{.bin}, GCC still compiled it as an ELF executable by default. The ELF file begins with the ASCII signature \texttt{0x7F 'E' 'L' 'F'}. The Abyss OS loader copied this file to \texttt{0x400000} and jumped to it, forcing the CPU to attempt executing the ELF headers as raw x86 machine code, resulting in an invalid opcode crash. Furthermore, an earlier bug involving \texttt{int 0x80} from Ring 0 caused stack misalignment (missing \texttt{ESP} and \texttt{SS}), confusing the \texttt{iret} instruction. \\
\textbf{Solution:} Replaced the syscalls with a direct C++ function pointer passed to the program. Crucially, the compiler command was updated to pass \texttt{--oformat,binary} to the linker, stripping the ELF metadata and producing a pure flat binary.

\subsection*{Phase 14--15: Multitasking \& Context Switching}

\subsubsection*{Bug 1: Dereferencing the Stack Pointer Instead of Using Its Address}
\textbf{Symptom:} Instant Page Fault on first context switch.\\
\textbf{Root Cause:} In \texttt{create\_task}, the stack pointer was saved incorrectly:
\begin{lstlisting}[language=C++]
// WRONG: Reads the VALUE at the top of the stack (which is 0, a dummy register)
new_task->esp = *stack;
new_task->ebp = *stack;

// CORRECT: Saves the MEMORY ADDRESS of the stack top
new_task->esp = (uint32_t)stack;
new_task->ebp = (uint32_t)stack;
\end{lstlisting}
Setting \texttt{esp = 0} pointed the new task's stack at the null pointer (address 0), causing an immediate fault when the CPU tried to \texttt{popa}.

\subsubsection*{Bug 2: Wrong \texttt{extern} Declaration in Assembly}
\textbf{Symptom:} NASM assembler error: \texttt{symbol `current\_task' not defined}.\\
\textbf{Root Cause:} The \texttt{extern} line in \texttt{boot/interrupts.asm} was accidentally typed as \texttt{extern create\_task} instead of \texttt{extern current\_task}. The variable the assembly code needed was never declared visible to the assembler.

\subsubsection*{Bug 3: Loading the Old Task's Stack Instead of the New Task's}
\textbf{Symptom:} The context switch ran but immediately jumped back to the same task -- no actual switching occurred.\\
\textbf{Root Cause:} In the assembly switch, \texttt{eax} holds the old task and \texttt{ebx} holds the new task. The stack restore used the wrong register:
\begin{lstlisting}[language={[x86masm]Assembler}]
; WRONG: Restores the OLD task's stack (no switch happens!)
mov esp, [eax+4]
mov ebp, [eax+8]

; CORRECT: Restores the NEW task's stack
mov esp, [ebx+4]
mov ebp, [ebx+8]
\end{lstlisting}

\subsubsection*{Bug 4: Incorrect IRET Stack Frame for New Tasks}
\textbf{Symptom:} Page Fault when switching to a brand-new task for the first time (never happened for already-running tasks).\\
\textbf{Root Cause:} When a hardware interrupt fires, the CPU automatically pushes \texttt{EFLAGS}, \texttt{CS}, and \texttt{EIP} onto the stack before our assembly code runs. The handler ends with \texttt{iret} which pops all three. But the fake stack we created for new tasks only faked the \texttt{pusha} registers, not the IRET frame. The solution was to push three extra values before the registers:
\begin{lstlisting}[language=C++]
*(--stack) = 0x202; // EFLAGS (interrupts enabled)
*(--stack) = 0x08;  // CS (Kernel Code Segment)
*(--stack) = (uint32_t)main_function; // EIP
// Then 8 dummy zeros for popa
for (int i = 0; i < 8; i++) *(--stack) = 0;
\end{lstlisting}

% ─────────────────────────────────────────────
\subsection*{Phase 16: Virtual File System}

\subsubsection*{Bug 5: Missing Type on Static Variable}
\textbf{Symptom:} Compiler error: \texttt{storage class specified for parameter}.\\
\textbf{Root Cause:} The \texttt{file\_count} variable was declared without a type:
\begin{lstlisting}[language=C++]
static file_count = 0;   // WRONG: no type
static uint32_t file_count = 0; // CORRECT
\end{lstlisting}

\subsubsection*{Bug 6: Using Wrong Field Name}
\textbf{Symptom:} Compiler error: \texttt{fs\_node has no member named 'size'}.\\
\textbf{Root Cause:} The field in \texttt{fs\_node} was named \texttt{length}, not \texttt{size}. The fix was simply \texttt{nodes[i].length = file\_size}.

\subsubsection*{Bug 7: Casting to a Function Pointer Instead of a Struct Pointer}
\textbf{Symptom:} Compiler error on the pointer arithmetic that advances to the next TAR block.\\
\textbf{Root Cause:} A typo cast the pointer to the wrong type:
\begin{lstlisting}[language=C++]
hdr = (tarfs_read*)((char*)hdr + 512 + blocks * 512); // WRONG
hdr = (tar_header*)((char*)hdr + 512 + blocks * 512); // CORRECT
\end{lstlisting}

\subsubsection*{Bug 8: Backwards \texttt{strcmp} Logic}
\textbf{Symptom:} \texttt{initrd\_findfile("hello.txt")} always returned null even though the file was in the archive.\\
\textbf{Root Cause:} The custom \texttt{strcmp} in this kernel returns \texttt{true} when strings are equal (opposite of the C standard which returns 0). The check was written with the wrong expectation:
\begin{lstlisting}[language=C++]
if (strcmp(nodes[i].name, file) == 0) // WRONG: true != 0
if (strcmp(nodes[i].name, file))      // CORRECT: true is truthy
\end{lstlisting}

\subsubsection*{Bug 9: Missing \texttt{\#include} in Header File}
\textbf{Symptom:} Compiler error: \texttt{fs\_node does not name a type} inside \texttt{initrd.h}.\\
\textbf{Root Cause:} \texttt{initrd.h} used the \texttt{fs\_node} type without including the header that defined it. The fix was adding \texttt{\#include "vfs.h"} at the top of \texttt{initrd.h}.

% ─────────────────────────────────────────────
\subsection*{Phase 17: VESA Graphics}

\subsubsection*{Bug 10: Framebuffer in Unmapped Virtual Memory (Silent Page Fault)}
\textbf{Symptom:} Completely black screen. No pixels drawn. No error visible.\\
\textbf{Root Cause:} This was the most subtle bug of the entire project. QEMU places the VESA linear framebuffer at physical address \texttt{0xFD000000}. However, the original \texttt{init\_paging()} only created one page table covering the first 4MB of RAM (\texttt{0x00000000}--\texttt{0x003FFFFF}). Every call to \texttt{put\_pixel} triggered a Page Fault because \texttt{0xFD000000} had no valid page table entry. The page fault handler ran invisibly (VGA text mode is ignored in graphics mode) and the kernel halted silently.\\
\textbf{Solution:} Enable PSE (Page Size Extension) in \texttt{CR4} and replace the single page table with 1024 identity-mapped 4MB Big Page entries, covering the entire 4GB address space:
\begin{lstlisting}[language=C++]
void init_paging() {
    // 0x83 = Present (1) + RW (2) + Page Size 4MB (0x80)
    for (int i = 0; i < 1024; i++)
        page_directory.tablesPhysical[i] = (i << 22) | 0x83;
    load_page_directory(page_directory.tablesPhysical);
    enable_paging();
}
\end{lstlisting}

\subsubsection*{Bug 11: Black and White Stripes Instead of Colored Rectangle}
\textbf{Symptom:} The screen displayed uniform alternating black and white vertical stripes instead of a solid colored rectangle.\\
\textbf{Root Cause:} The \texttt{framebuffer\_addr} field in \texttt{multiboot\_info} is declared as \texttt{uint64\_t} (8 bytes). Passing it directly to \texttt{vesa\_init(uint32\_t addr, ...)} without an explicit cast caused the compiler to push two 32-bit words onto the stack. Inside the function, the parameters \texttt{width}, \texttt{height}, and \texttt{pitch} were all read from the wrong stack offsets. The \texttt{pitch} value received was completely wrong, making every row of pixels map to the wrong screen address.\\
\textbf{Solution:} Explicit cast at the call site:
\begin{lstlisting}[language=C++]
vesa_init((uint32_t)mbd->framebuffer_addr,
           mbd->framebuffer_width,
           mbd->framebuffer_height,
           mbd->framebuffer_pitch);
\end{lstlisting}

\subsubsection*{Bug 12: Grayscale Fade (Implicit Integer Promotion)}
\textbf{Symptom:} When rendering a gradient from Crimson to Dark Purple, the screen instead drew a bizarre grayscale fade from Black to White.\\
\textbf{Root Cause:} The formula \texttt{uint8\_t r = r1 + (((int)r2 - r1) * i) / h;} suffered from \textbf{Implicit Integer Promotion}. \texttt{i} and \texttt{h} were defined as \texttt{uint32\_t} (unsigned integers). If \texttt{r2} was darker than \texttt{r1}, the result of \texttt{r2 - r1} was a negative number. When C++ attempted to multiply this signed negative number by the unsigned \texttt{i}, it automatically converted the negative number into an unsigned format, wrapping it around into a massive positive number (e.g., 4,294,967,280). This overflowed the math completely.\\
\textbf{Solution:} Explicitly cast the loop variables to signed integers before the math:
\begin{lstlisting}[language=C++]
uint8_t r = r1 + (((int)r2 - (int)r1) * (int)i) / (int)h;
\end{lstlisting}

\subsubsection*{Bug 13: The Infinite Text Loop}
\textbf{Symptom:} The OS locked up completely when trying to print a string, freezing the Virtual Machine.\\
\textbf{Root Cause:} The \texttt{draw\_string} function contained a \texttt{while(*str != '\textbackslash 0')} loop, but we forgot to increment the character pointer (\texttt{str++}) inside the loop. The CPU endlessly drew the exact same first letter on top of itself forever.\\
\textbf{Solution:} Add \texttt{str++;} to the bottom of the loop.

\subsubsection*{Bug 14: Text Drawing Off the Edge of the Screen}
\textbf{Symptom:} Multi-line strings were drawn as one massive continuous horizontal line that shot completely off the right edge of the monitor.\\
\textbf{Root Cause:} The text renderer had no logic to understand the newline character (\texttt{\textbackslash n}). It just treated it as an unprintable space, incremented \texttt{X}, and kept drawing to the right.\\
\textbf{Solution:} Added a conditional check inside \texttt{draw\_string} to simulate a Carriage Return (reset X to starting margin) and a Line Feed (push Y down by the font height):
\begin{lstlisting}[language=C++]
if (*str == '\n') {
    x = start_x;
    y += (8 * scale);
}
\end{lstlisting}

\subsubsection*{Bug 15: Font Array 'Not Declared'}
\textbf{Symptom:} Compilation failed with the error: \texttt{error: 'font8x8' was not declared in this scope}.\\
\textbf{Root Cause:} The standard open-source header file we downloaded named the array \texttt{font8x8\_basic}, but our \texttt{vesa.cpp} code attempted to reference it purely as \texttt{font8x8}.\\
\textbf{Solution:} Changed the reference in \texttt{draw\_char} to match the header exactly: \texttt{font8x8\_basic[(uint8\_t)c][row]}.

\subsubsection*{Bug 16: The Vertical Garbage Column (Missing Increment)}
\textbf{Symptom:} The OS booted, but the first letter of a large message was replaced with a glitchy vertical line of garbled text running down the left side of the screen.\\
\textbf{Root Cause:} In the newly refactored \texttt{vga\_print\_char}, we called \texttt{draw\_char(c, term\_x, term\_y, ...)} but forgot to increment \texttt{term\_x} afterwards. Every subsequent character printed by the OS was drawn at the exact same X and Y coordinate, stacking perfectly on top of each other and forming a solid glitch block.\\
\textbf{Solution:} Added \texttt{term\_x += 8 * FONT\_SCALE;} immediately after drawing the character.

\subsubsection*{Bug 17: Mixing Coordinate Spaces}
\textbf{Symptom:} Text from \texttt{vga\_print} was drawing directly over the top of text printed by \texttt{draw\_string}, completely ignoring formatting.\\
\textbf{Root Cause:} In \texttt{kernel\_main}, we called \texttt{draw\_string(..., 0, 0)} (which manages its own local X and Y variables) followed immediately by \texttt{vga\_print(...)} (which relies on the global \texttt{term\_x} and \texttt{term\_y} variables in \texttt{vga.cpp}). Because \texttt{draw\_string} did not update the global trackers, \texttt{vga\_print} assumed the cursor was still at `0, 0` and printed right over the previous string.\\
\textbf{Solution:} Replaced the manual \texttt{draw\_string} call in \texttt{kernel\_main} with standard \texttt{vga\_print} calls so all text shares the same global cursor context.

\subsubsection*{Bug 18: The Black Hole Terminal (Missing Reset)}
\textbf{Symptom:} The terminal scrolled down normally, but once it hit the bottom of the screen, the screen flashed black and no more text was ever seen again.\\
\textbf{Root Cause:} The \texttt{scroll()} function successfully detected that \texttt{term\_y} was off-screen and triggered \texttt{vga\_clear()}. \texttt{vga\_clear()} successfully drew a black rectangle to clear the screen, but it forgot to reset \texttt{term\_x} and \texttt{term\_y} back to 0. Consequently, \texttt{term\_y} remained off-screen indefinitely, meaning every subsequent character immediately triggered another invisible clear.\\
\textbf{Solution:} Added \texttt{term\_x = TERM\_START\_X;} and \texttt{term\_y = TERM\_START\_Y;} to the bottom of \texttt{vga\_clear()}.

\subsubsection*{Bug 19: The Top-Edge Magnetic Pinning}
\textbf{Symptom:} The mouse cursor became permanently pinned to the very top edge of the screen (\texttt{Y = 0}) and refused to move down, although X-axis movement worked fine.\\
\textbf{Root Cause:} The PS/2 packet sends movement deltas as 8-bit two's complement integers. We casted the raw byte to a \texttt{char}, but our C++ compiler configuration defaulted to treating \texttt{char} as an \textit{unsigned} 8-bit integer. When the mouse sent a negative delta (e.g., \texttt{-1} or \texttt{0xFF}), the compiler interpreted it as \texttt{+255}. Subtracting \texttt{255} from \texttt{mouse\_y} caused it to plunge instantly into the negatives, where our bounds-checking clamped it permanently at \texttt{0}.\\
\textbf{Solution:} Explicitly cast the delta bytes to \texttt{int8\_t} to force standard sign-extension, allowing negative movements to be properly interpreted.

\subsubsection*{Bug 20: Missing Interrupt Framework Abstraction}
\textbf{Symptom:} Compilation failed with \texttt{isr.h: No such file or directory} when attempting to hook IRQ 12 using a generic \texttt{register\_interrupt\_handler} function.\\
\textbf{Root Cause:} We attempted to use a generic interrupt registration abstraction commonly found in standard OS tutorials. However, Abyss OS utilizes a highly customized, direct-to-assembly IDT architecture for speed, meaning no such registry existed.\\
\textbf{Solution:} Dropped the generic approach and directly wired \texttt{isr44} into \texttt{boot/interrupts.asm}, added the gate to \texttt{kernel/idt.cpp}, and exposed \texttt{extern "C" void isr44\_handler()} natively in \texttt{mouse.cpp}.

\newpage

\section{Phase 20: The Graphical User Interface \& Double Buffering}
{\small\textit{\textbf{Completed:} August 29, 2026}}\par\medskip
\subsection{The Double Buffering Architecture}
When the mouse driver was first implemented, it wrote directly to the VESA video memory (\texttt{framebuffer}). This meant the driver had to manually save the background pixels before drawing the cursor, and restore them when the cursor moved. However, because the background was a complex gradient, overlapping movements caused pixel corruption.
To solve this, we implemented Double Buffering. A 3.1 MB \texttt{backbuffer} was allocated on the heap. This exposed a critical memory bug: our initial bump allocator (\texttt{kmalloc}) was hard-capped at 4MB, and was never initialized in \texttt{kernel\_main}. We expanded the heap limit to 32MB and correctly initialized the memory subsystem. 
All drawing functions (\texttt{put\_pixel}, \texttt{draw\_rect}) were refactored to draw to the backbuffer in software.

\subsection{The Compositor Loop}
A new \texttt{ui\_render()} function was created to act as a Compositor. Instead of the kernel sleeping permanently, the UI loop clears the backbuffer, draws the gradient, draws any active GUI Windows, draws the mouse cursor on top, and finally swaps the buffers to the screen using a blistering-fast inline assembly instruction: \texttt{rep movsl}.

\section{Phase 21: The Interactive Desktop}
{\small\textit{\textbf{Completed:} August 29, 2026}}\par\medskip
\subsection{Click Detection and Z-Order Promotion}
The PS/2 Mouse packet's first byte contains the button states. We extracted Bit 0 to track \texttt{mouse\_left\_pressed}. We introduced a \texttt{Window} struct to hold window coordinates, dimensions, and drag states.
In the compositor loop, a reverse-iteration (Z-Order) collision detection was implemented. If the user clicks the title bar of a window, it is pulled from its current Z-index, shifted to the end of the rendering array (promoting it to the foreground), and its offset is recorded to allow smooth dragging without snapping to the cursor center.

\newpage


\section{Phase 22: The ATA Hard Drive Driver}
{\small\textit{\textbf{Completed:} September 2, 2026}}\par\medskip

\subsection{Theoretical Background}
Before writing the driver, we must understand the core technologies governing physical hard drives in legacy systems:
\begin{itemize}
    \item \textbf{ATA (Advanced Technology Attachment):} The standard physical interface used to connect storage devices (Hard Drives, CD-ROMs) to the PC motherboard.
    \item \textbf{PIO (Programmed Input/Output):} A data transmission method where the CPU itself must manually execute an instruction (like \texttt{inw} or \texttt{outw}) to transfer every single byte of data between the hardware and RAM. While slower than modern DMA (Direct Memory Access), it is the most reliable baseline for storage drivers.
    \item \textbf{LBA (Logical Block Addressing):} Instead of addressing data physically by Cylinder, Head, and Sector (CHS), LBA abstracts the drive into a flat array of 512-byte blocks. We use LBA28, which uses a 28-bit address, allowing the OS to access up to 128 GB of storage.
\end{itemize}

\subsection{The Objective}
Our objective in this phase was to break free from the read-only constraints of the Initial Ramdisk (Initrd). We needed to:
\begin{enumerate}
    \item Generate a virtual 64MB raw image file (\texttt{disk.img}) and attach it to QEMU as the Primary Master drive.
    \item Implement hardware-accelerated 16-bit port I/O to read and write entire 512-byte sectors.
    \item Write an ATA PIO driver to permanently etch a message onto the magnetic disk, and verify it via a hexadecimal dump on the host machine.
\end{enumerate}

\subsection{Implementation \& Bug Resolution}
Because hard drives process data in 512-byte sectors (256 16-bit words), relying on our standard 8-bit \texttt{inb} port reader would be highly inefficient. We expanded \texttt{io.h} to include the x86 \texttt{rep insw} and \texttt{rep outsw} inline assembly instructions, allowing the CPU to blast 16-bit words directly between the IDE controller and system memory at maximum hardware speeds.

The ATA driver sends the 28-bit LBA sector address across ports \texttt{0x1F3} through \texttt{0x1F6}, and sends the READ (\texttt{0x20}) or WRITE (\texttt{0x30}) command to \texttt{0x1F7}. 

During testing, a critical interrupt conflict was discovered: when the ATA controller finishes a write command, it fires IRQ 14 (Interrupt 46). Since our IDT lacked an \texttt{isr46} handler, the OS crashed immediately after the disk write. Because our driver uses polling (waiting for \texttt{BSY} and \texttt{DRQ} flags rather than using interrupts), we elegantly resolved the crash by setting the \texttt{nIEN} (No Interrupt Enable) bit in the Device Control Register (\texttt{0x3F6}), instructing the hard drive to remain silent after finishing commands. 

\newpage

\section{Phase 23: The FAT32 File System}
{\small\textit{\textbf{Completed:} September 6, 2026}}\par\medskip

\subsection{Theoretical Background}
Reading raw bytes from a hard drive using LBA sectors is only half the battle. To provide a usable operating system, we must implement a standard File System. We chose \textbf{FAT32} (File Allocation Table 32-bit), originally introduced in Windows 95 OSR2, because it is the most widely supported legacy file system and remains the standard for USB flash drives today.

A FAT32 volume consists of several key regions:
\begin{itemize}
    \item \textbf{VBR (Volume Boot Record):} Located at Sector 0 of the partition. It contains crucial metadata about the file system, including the number of sectors per cluster and the location of the Root Directory.
    \item \textbf{FAT (File Allocation Table):} A massive array of 32-bit integers that act as a linked list. It tracks which clusters belong to which files, allowing files to be fragmented across the physical disk.
    \item \textbf{Data Region:} The actual sectors where file contents and directories are stored, organized into "Clusters" (blocks of contiguous sectors).
\end{itemize}

\subsection{The Objective}
The goal of this phase was to teach Abyss OS to parse the FAT32 VBR and read the Root Directory. We needed to format our virtual \texttt{disk.img} as a FAT32 volume, inject a test file from the Linux host machine, and successfully print the name of that file to the Abyss OS screen.

\subsection{Implementation \& Structure Alignment}
We utilized the Linux \texttt{mtools} utility (\texttt{mformat} and \texttt{mcopy}) to format \texttt{disk.img} and inject \texttt{HELLO.TXT} without requiring root privileges. 

To parse the VBR, we had to ensure our C++ \texttt{FAT32\_VBR} structure perfectly matched the physical byte layout on the disk. Because C++ compilers often optimize structures by injecting invisible "padding bytes" to align variables to 4-byte memory boundaries, reading 512 raw bytes directly into a standard struct would result in misaligned garbage data. 

To solve this, we used compiler directives to force exact physical byte alignment:
\begin{lstlisting}[language=C++]
#pragma pack(push, 1) // Instructs compiler to pack variables tightly (1-byte alignment)

struct FAT32_VBR { ... };
struct FAT32_DirectoryEntry { ... };

#pragma pack(pop)     // Restores normal compiler alignment optimizations
\end{lstlisting}

By casting the 512 bytes from Sector 0 into the \texttt{FAT32\_VBR} pointer, the kernel successfully dynamically calculated the LBA of the Root Directory cluster. It then loaded the Root Directory into memory, cast the buffer into an array of 32-byte \texttt{FAT32\_DirectoryEntry} structures, filtered out deleted or hidden files, and successfully rendered the \texttt{HELLO.TXT} filename to the GUI.

\newpage

\section{Phase 24: Reading FAT32 Clusters}
{\small\textit{\textbf{Completed:} September 20, 2026}}\par\medskip

\subsection{Theoretical Background}
In Phase 23, we parsed the Root Directory and found the 32-byte metadata entry for \texttt{HELLO.TXT}. However, a FAT32 directory entry only stores the \textit{starting cluster} of a file. It does not tell you where the rest of the file is located! 

Because operating systems constantly create, delete, and resize files, files are rarely stored in contiguous blocks on the hard drive. Instead, they are fragmented. To reconstruct a fragmented file, FAT32 uses the \textbf{File Allocation Table (FAT)}. The FAT is a massive array of 32-bit integers that acts as a linked list mapping out the physical disk.

\subsection{The Objective}
The goal of this phase was to write a cluster traversal algorithm that could query the FAT table, follow a file's linked list chain across the physical disk, extract the raw data sectors into RAM, and reconstruct the complete \texttt{HELLO.TXT} string inside the Abyss OS environment.

\subsection{Implementation: Traversing the Chain}
Because the FAT table itself is stored on the hard drive, looking up a cluster requires a disk read. We implemented \texttt{fat32\_get\_next\_cluster(current\_cluster)}, which performs the following calculations:
\begin{enumerate}
    \item \textbf{Sector Offset:} Since 1 sector is 512 bytes, and each FAT entry is 4 bytes (32-bit), there are exactly 128 entries per sector. We divide the cluster ID by 128 to find which physical sector of the FAT table to read.
    \item \textbf{Byte Offset:} We take the cluster ID modulo 128 and multiply by 4 to find the exact byte offset of the entry inside that sector.
\end{enumerate}

After reading the FAT sector via the ATA driver, the 32-bit integer is extracted. Because FAT32 only uses 28 bits, the top 4 bits are masked out (\texttt{0x0FFFFFFF}). This value is the next cluster in the file's chain. A value of \texttt{0x0FFFFFF8} or higher signifies the End-of-File (EOF).

We successfully integrated this logic into \texttt{fat32\_read\_file}, which allocates a buffer, reads the data clusters, and dynamically follows the FAT chain until EOF is reached. Our testing proved flawless: Abyss OS allocated 2KB of RAM, read the fragmented clusters, and rendered the internal text of \texttt{HELLO.TXT} directly to the GUI!

\newpage

\section{Phase 25: The Graphical File Explorer}
{\small\textit{\textbf{Completed:} September 20, 2026}}\par\medskip

\subsection{Theoretical Background}
In modern operating systems, the file system backend and the graphical user interface (GUI) frontend are completely separate subsystems. The GUI relies on a robust API to query the file system for directory entries and data. Up until this point, Abyss OS had only tested the FAT32 driver by printing text directly to the screen via a blocking \texttt{while(1)} loop.

\subsection{The Objective}
The goal of this phase was to construct Abyss OS's first native, data-driven graphical application. We needed to integrate the FAT32 driver directly into the window manager, allowing a "File Explorer" window to dynamically query the hard drive and render its contents as graphical elements on the screen.

\subsection{Implementation: Building the App}
To decouple the backend from the frontend, we introduced a new API function in \texttt{kernel/fat32.cpp}: \texttt{fat32\_get\_root\_entries()}. This function queries the FAT32 root directory on the physical disk and populates an array of \texttt{FAT32\_DirectoryEntry} structures, abstracting away the low-level ATA sector reads from the GUI.

In the window manager (\texttt{kernel/ui.cpp}), we expanded our global \texttt{windows} array to include a third application: "File Explorer". During the primary \texttt{ui\_render()} loop, if the renderer detects it is drawing the File Explorer window, it executes custom application logic:
\begin{enumerate}
    \item It clears the interior of the window with a solid white background.
    \item It calls the \texttt{fat32\_get\_root\_entries()} API to fetch the current files from the hard drive.
    \item It loops through each returned file, utilizing our primitive graphics functions to draw a blue icon (\texttt{draw\_rect}) and render the 8.3 filename (\texttt{draw\_string}) directly into the application's coordinate space.
\end{enumerate}

This successfully demonstrated a true OS architecture: a hardware driver providing abstracted data to a high-level graphical window manager!

\newpage

\chapter{Complete Source Code}


\subsection{Makefile}
\begin{lstlisting}[language={bash}]
CC = i686-elf-g++
ASM = nasm
LD = i686-elf-ld

CFLAGS = -ffreestanding -O2 -Wall -Wextra -fno-exceptions -fno-rtti
LDFLAGS = -T linker.ld -ffreestanding -O2 -nostdlib

all: abyss.iso

boot/boot.o: boot/boot.asm
	$(ASM) -f elf32 boot/boot.asm -o boot/boot.o

boot/gdt_flush.o: boot/gdt_flush.asm
	$(ASM) -f elf32 boot/gdt_flush.asm -o boot/gdt_flush.o

boot/idt_flush.o: boot/idt_flush.asm
	$(ASM) -f elf32 boot/idt_flush.asm -o boot/idt_flush.o

boot/interrupts.o: boot/interrupts.asm
	$(ASM) -f elf32 boot/interrupts.asm -o boot/interrupts.o

boot/paging.o: boot/paging.asm
	$(ASM) -f elf32 boot/paging.asm -o boot/paging.o

boot/usermode.o: boot/usermode.asm
	$(ASM) -f elf32 boot/usermode.asm -o boot/usermode.o

boot/task.o: boot/task.asm
	$(ASM) -f elf32 boot/task.asm -o boot/task.o

kernel/kernel.o: kernel/kernel.cpp
	$(CC) $(CFLAGS) -c kernel/kernel.cpp -o kernel/kernel.o

kernel/gdt.o: kernel/gdt.cpp
	$(CC) $(CFLAGS) -c kernel/gdt.cpp -o kernel/gdt.o

kernel/idt.o: kernel/idt.cpp
	$(CC) $(CFLAGS) -c kernel/idt.cpp -o kernel/idt.o

kernel/vga.o: kernel/vga.cpp
	$(CC) $(CFLAGS) -c kernel/vga.cpp -o kernel/vga.o

kernel/isr.o: kernel/isr.cpp
	$(CC) $(CFLAGS) -c kernel/isr.cpp -o kernel/isr.o

kernel/pic.o: kernel/pic.cpp
	$(CC) $(CFLAGS) -c kernel/pic.cpp -o kernel/pic.o

kernel/paging.o: kernel/paging.cpp	
	$(CC) $(CFLAGS) -c kernel/paging.cpp -o kernel/paging.o

kernel/string.o: kernel/string.cpp
	$(CC) $(CFLAGS) -c kernel/string.cpp -o kernel/string.o

kernel/tss.o: kernel/tss.cpp
	$(CC) $(CFLAGS) -c kernel/tss.cpp -o kernel/tss.o

kernel/memory.o: kernel/memory.cpp
	$(CC) $(CFLAGS) -c kernel/memory.cpp -o kernel/memory.o

kernel/timer.o: kernel/timer.cpp
	$(CC) $(CFLAGS) -c kernel/timer.cpp -o kernel/timer.o 

kernel/task.o: kernel/task.cpp
	$(CC) $(CFLAGS) -c kernel/task.cpp -o kernel/task.o

kernel/vfs.o: kernel/vfs.cpp
	$(CC) $(CFLAGS) -c kernel/vfs.cpp -o kernel/vfs.o 

kernel/initrd.o: kernel/initrd.cpp
	$(CC) $(CFLAGS) -c kernel/initrd.cpp -o kernel/initrd.o 

kernel/vesa.o: kernel/vesa.cpp
	$(CC) $(CFLAGS) -c kernel/vesa.cpp -o kernel/vesa.o

kernel/mouse.o: kernel/mouse.cpp
	$(CC) $(CFLAGS) -c kernel/mouse.cpp -o kernel/mouse.o 

kernel/ui.o: kernel/ui.cpp
	$(CC) $(CFLAGS) -c kernel/ui.cpp -o kernel/ui.o

kernel/ata.o: kernel/ata.cpp
	$(CC) $(CFLAGS) -c kernel/ata.cpp -o kernel/ata.o

kernel/fat32.o: kernel/fat32.cpp
	$(CC) $(CFLAGS) -c kernel/fat32.cpp -o kernel/fat32.o 

abyss.bin: boot/boot.o kernel/kernel.o boot/gdt_flush.o boot/idt_flush.o boot/interrupts.o boot/paging.o kernel/gdt.o kernel/idt.o kernel/vga.o kernel/isr.o kernel/pic.o kernel/paging.o kernel/string.o boot/usermode.o kernel/tss.o kernel/memory.o kernel/timer.o kernel/task.o boot/task.o kernel/vfs.o kernel/initrd.o kernel/vesa.o kernel/mouse.o kernel/ui.o kernel/ata.o kernel/fat32.o
	$(CC) $(LDFLAGS) -o abyss.bin boot/boot.o kernel/kernel.o boot/gdt_flush.o boot/idt_flush.o boot/interrupts.o boot/paging.o kernel/gdt.o kernel/idt.o kernel/vga.o kernel/isr.o kernel/pic.o kernel/paging.o kernel/string.o boot/usermode.o kernel/tss.o kernel/memory.o kernel/timer.o kernel/task.o boot/task.o kernel/vfs.o kernel/initrd.o kernel/vesa.o kernel/mouse.o kernel/ui.o kernel/ata.o kernel/fat32.o

abyss.iso: abyss.bin
	cp abyss.bin iso/boot/abyss.bin
	cd initrd && tar --format=ustar -cvf ../iso/boot/initrd.tar *
	grub-mkrescue -o abyss.iso iso

run: abyss.iso
# 	qemu-system-i386 -cdrom abyss.iso
	qemu-system-i386 -hda disk.img -cdrom abyss.iso -boot d -vga std

clean:
	rm -f boot/boot.o kernel/kernel.o abyss.bin abyss.iso iso/boot/abyss.bin boot/gdt_flush.o boot/idt_flush.o boot/interrupts.o boot/paging.o kernel/gdt.o kernel/idt.o kernel/vga.o kernel/isr.o kernel/pic.o kernel/paging.o kernel/string.o boot/usermode.o kernel/tss.o kernel/memory.o kernel/timer.o kernel/task.o boot/task.o kernel/vfs.o kernel/initrd.o kernel/vesa.o kernel/mouse.o kernel/ui.o kernel/ata.o kernel/fat32.o

\end{lstlisting}

\subsection{boot.asm}
\begin{lstlisting}[language={[x86masm]Assembler}]
MAGIC    equ 0x1BADB002
FLAGS    equ 0x04
CHECKSUM equ -(MAGIC + FLAGS)

section .multiboot
align 4
    dd MAGIC
    dd FLAGS
    dd CHECKSUM
    dd 0, 0, 0, 0, 0
    dd 0
    dd 1024
    dd 768
    dd 32

section .bss
align 16
stack_bottom:
    resb 16384      ; reserve 16 KB
stack_top:

section .text
global _start
extern kernel_main

_start:
    mov esp, stack_top
    push ebx
    call kernel_main
    cli
.hang:
    hlt
    jmp .hang

\end{lstlisting}

\subsection{gdt\_flush.asm}
\begin{lstlisting}[language={[x86masm]Assembler}]
global gdt_flush
section .text
gdt_flush:
    mov eax, [esp+4]
    lgdt [eax]
    mov ax, 0x10
    mov ds, ax
    mov es, ax
    mov fs, ax
    mov gs, ax
    mov ss, ax
    jmp 0x08:.flush
    .flush:
        ret
\end{lstlisting}

\subsection{idt\_flush.asm}
\begin{lstlisting}[language={[x86masm]Assembler}]
global idt_flush
section .text
idt_flush:
    mov eax, [esp+4]
    lidt [eax]
    ret
\end{lstlisting}

\subsection{interrupts.asm}
\begin{lstlisting}[language={[x86masm]Assembler}]
section .text

; --- MACROS TO GENERATE 32 EXCEPTION HANDLERS ---
%macro ISR_NOERRCODE 1
  global isr%1
  isr%1:
    push 0   
    push %1 
    jmp isr_common_stub
%endmacro

%macro ISR_ERRCODE 1
  global isr%1
  isr%1:
    push %1
    jmp isr_common_stub
%endmacro

; Generate ISR 0 to 31
ISR_NOERRCODE 0
ISR_NOERRCODE 1
ISR_NOERRCODE 2
ISR_NOERRCODE 3
ISR_NOERRCODE 4
ISR_NOERRCODE 5
ISR_NOERRCODE 6
ISR_NOERRCODE 7
ISR_ERRCODE   8
ISR_NOERRCODE 9
ISR_ERRCODE   10
ISR_ERRCODE   11
ISR_ERRCODE   12
ISR_ERRCODE   13
ISR_ERRCODE   14
ISR_NOERRCODE 15
ISR_NOERRCODE 16
ISR_ERRCODE   17
ISR_NOERRCODE 18
ISR_NOERRCODE 19
ISR_NOERRCODE 20
ISR_NOERRCODE 21
ISR_NOERRCODE 22
ISR_NOERRCODE 23
ISR_NOERRCODE 24
ISR_NOERRCODE 25
ISR_NOERRCODE 26
ISR_NOERRCODE 27
ISR_NOERRCODE 28
ISR_NOERRCODE 29
ISR_NOERRCODE 30
ISR_NOERRCODE 31
ISR_NOERRCODE 128

; Send EOI (IRQ 12) to the PIC!
; --- THE COMMON C++ CALLER ---
extern fault_handler

isr_common_stub:
    pusha        ; Save all registers
    push esp     ; Pass a pointer to the stack to our C++ function!
    call fault_handler
    add esp, 4   ; Clean up the pointer we just pushed
    popa         ; Restore registers
    add esp, 8   ; Clean up the pushed error code and pushed ISR number
    iret

; --- Interrupt 32 (Timer IRQ0) ---
global isr32
extern isr32_handler
extern current_task
isr32:
    pusha
    call isr32_handler

    mov eax, [current_task]
    mov [eax+4], esp
    mov [eax+8], ebp
    
    mov ebx, [eax+16]
    mov [current_task], ebx

    mov esp, [ebx+4]
    mov ebp, [ebx+8]

    popa
    iret

; --- Interrupt 33 (Keyboard IRQ1) ---
global isr33
extern isr33_handler
isr33:
    pusha
    call isr33_handler
    popa
    iret

global isr44
extern isr44_handler
isr44:
    pusha
    call isr44_handler
    popa
    iret

\end{lstlisting}

\subsection{paging.asm}
\begin{lstlisting}[language={[x86masm]Assembler}]
global load_page_directory
global enable_paging

section .text

; void load_page_directory(uint32_t* pointer)
load_page_directory:
    push ebp
    mov ebp, esp
    
    mov eax, [ebp+8]   ; Grab the pointer we passed from C++
    mov cr3, eax       ; Load it into the magical CR3 register
    
    mov esp, ebp
    pop ebp
    ret

; void enable_paging()
enable_paging:
    push ebp
    mov ebp, esp
    ; Enable PSE (Page Size Extension) so we can use 4MB pages
    mov eax, cr4
    or  eax, 0x10 
    mov cr4, eax

    mov eax, cr0
    or eax, 0x80000000 ; Flip the 31st bit (the Paging Bit) to 1
    mov cr0, eax       ; Push it back into the CPU! (Paging turns ON right here)
    
    mov esp, ebp
    pop ebp
    ret

\end{lstlisting}

\subsection{task.asm}
\begin{lstlisting}[language={[x86masm]Assembler}]
section .text

global perform_task_switch
extern current_task

perform_task_switch:
    ; 1. Save the CURRENT task's state onto its own stack!
    pushf           ; Push EFLAGS register
    pusha           ; Push all 8 general purpose registers
    
    ; 2. Get the pointer to the current_task struct in memory
    mov eax, [current_task]
    
    ; 3. Save the current stack pointers into the struct!
    ; (id is at offset 0. esp is at offset 4. ebp is at offset 8)
    mov [eax+4], esp
    mov [eax+8], ebp
    
    ; 4. Move to the next task in the linked list
    ; (next is at offset 16, assuming eip is still at offset 12 in your struct!)
    mov ebx, [eax+16]
    mov [current_task], ebx
    
    ; 5. THE SWITCH! Load the NEW task's stack pointers into the CPU!
    mov esp, [ebx+4]
    mov ebp, [ebx+8]
    
    ; 6. Restore the NEW task's state from its own stack!
    popa            ; Pop all 8 general purpose registers
    popf            ; Pop EFLAGS register
    
    ; 7. Because C++ called this function, 'ret' automatically pops the 
    ;    Instruction Pointer (EIP) off the stack and jumps to it!
    ret

\end{lstlisting}

\subsection{usermode.asm}
\begin{lstlisting}[language={[x86masm]Assembler}]
global jump_usermode
extern test_user_function ; This is a C++ function we will write next!

section .text
jump_usermode:
    cli  ; Turn off hardware interrupts so we don't crash while building the fake stack

    ; 1. Set the Data Segment registers to our new Ring 3 Data GDT Entry
    ; (Our entry is 0x20. We add 3 to tell the CPU it is Ring 3 -> 0x23)
    mov ax, 0x23
    mov ds, ax
    mov es, ax
    mov fs, ax
    mov gs, ax

    ; 2. Build the fake interrupt stack!
    push 0x23                ; SS (Stack Segment)
    push esp                 ; ESP (We'll just reuse the current stack)
    pushf                    ; EFLAGS (Push current flags)
    push 0x1B                ; CS (Ring 3 Code Entry is 0x18 + 3 = 0x1B)
    push test_user_function  ; EIP (The C++ function we want to run in Ring 3!)

    ; 3. Execute the fake return!
    iret

\end{lstlisting}

\subsection{hello.txt}
\begin{lstlisting}[language={bash}]
Greetings from the Abyss' Virtual file system, mortal!!

\end{lstlisting}

\subsection{help.txt}
\begin{lstlisting}[language={bash}]
"These are the helping commands mortal :\n 1. clear\n 2. hail\n 3. alloc/malloc\n 4. uptime\n 5. cat <filename>\n 6. help\n 7. ls\n"

\end{lstlisting}

\subsection{ata.cpp}
\begin{lstlisting}[language={C++}]
#include "ata.h"
#include "io.h"

#define ATA_DATA        0x1F0
#define ATA_SECTOR_CNT  0x1F2
#define ATA_LBA_LOW     0x1F3
#define ATA_LBA_MID     0x1F4
#define ATA_LBA_HIGH    0x1F5
#define ATA_DRIVE       0x1F6
#define ATA_COMMAND     0x1F7
#define ATA_STATUS      0x1F7

#define ATA_CMD_READ_PIO  0x20
#define ATA_CMD_WRITE_PIO 0x30

#define ATA_SR_BSY     0x80    
#define ATA_SR_DRQ     0x08 

void ata_wait_bsy(){
    while(inb(ATA_STATUS) & ATA_SR_BSY);
}

void ata_wait_drq(){
    while(!(inb(ATA_STATUS) & ATA_SR_DRQ));
}

void ata_init(){
    outb(0x3F6, 0x02);
}

void ata_read_sector(uint32_t lba, uint8_t* buf){
    ata_wait_bsy();
    outb(ATA_DRIVE, 0xE0 | ((lba >> 24) & 0x0F));
    outb(ATA_SECTOR_CNT, 1);
    outb(ATA_LBA_LOW, (uint8_t)lba);
    outb(ATA_LBA_MID, (uint8_t)(lba >> 8));
    outb(ATA_LBA_HIGH, (uint8_t)(lba >> 16));

    outb(ATA_COMMAND, ATA_CMD_READ_PIO);
    ata_wait_bsy();
    ata_wait_drq();

    insw(ATA_DATA, buf, 256);
}

void ata_write_sector(uint32_t lba, uint8_t* buffer) {
    ata_wait_bsy();
    
    outb(ATA_DRIVE, 0xE0 | ((lba >> 24) & 0x0F));
    outb(ATA_SECTOR_CNT, 1);
    outb(ATA_LBA_LOW, (uint8_t)lba);
    outb(ATA_LBA_MID, (uint8_t)(lba >> 8));
    outb(ATA_LBA_HIGH, (uint8_t)(lba >> 16));
    
    outb(ATA_COMMAND, ATA_CMD_WRITE_PIO);
    
    ata_wait_bsy();
    ata_wait_drq();
    
    outsw(ATA_DATA, buffer, 256);
    ata_wait_bsy();
    outb(ATA_COMMAND, 0xE7); 
    ata_wait_bsy();
}
\end{lstlisting}

\subsection{ata.h}
\begin{lstlisting}[language={C++}]
#pragma once 
#include <stdint.h>

void ata_init();
void ata_read_sector(uint32_t lba, uint8_t* buf);
void ata_write_sector(uint32_t lba, uint8_t* buf);
\end{lstlisting}

\subsection{fat32.cpp}
\begin{lstlisting}[language={C++}]
#include "fat32.h"
#include "ata.h"
#include "vga.h"

uint32_t fat_start_sector = 0;
uint32_t data_start_sector = 0;
uint32_t root_dir_cluster = 0;
uint8_t  sectors_per_cluster = 0;

// Converts a FAT32 Cluster Number into a physical ATA LBA Sector Number
uint32_t cluster_to_lba(uint32_t cluster) {
    return data_start_sector + ((cluster - 2) * sectors_per_cluster);
}
bool str_match(const char* str1, const char* str2, int len) {
    for (int i = 0; i < len; i++) {
        if (str1[i] != str2[i]) return false;
    }
    return true;
}

void fat32_init() {
    uint8_t boot_sector[512];
    ata_read_sector(0, boot_sector); // Read LBA 0
    
    FAT32_VBR* vbr = (FAT32_VBR*)boot_sector;
    
    // Check if it's actually FAT32
    if (vbr->boot_signature != 0xAA55) {
        vga_print("Error: Disk is not FAT32 formatted!\n");
        return;
    }
    
    sectors_per_cluster = vbr->sectors_per_cluster;
    root_dir_cluster = vbr->root_cluster;
    
    // Calculate memory regions
    fat_start_sector = vbr->reserved_sectors;
    uint32_t fat_size = vbr->sectors_per_fat_32;
    data_start_sector = fat_start_sector + (vbr->fat_count * fat_size);
    
    vga_print("FAT32 Init Success! Sys ID: ");
    for(int i=0; i<8; i++) vga_print_char(vbr->sys_identifier[i]);
    vga_print("\n");
}

void fat32_list_root() {
    uint8_t root_buffer[512];
    uint32_t root_lba = cluster_to_lba(root_dir_cluster);
    
    // Read the first sector of the root directory
    ata_read_sector(root_lba, root_buffer);
    
    FAT32_DirectoryEntry* entries = (FAT32_DirectoryEntry*)root_buffer;
    
    // A 512-byte sector holds exactly sixteen 32-byte directory entries (512 / 32 = 16)
    for (int i = 0; i < 16; i++) {
        FAT32_DirectoryEntry* entry = &entries[i];
        
        // 0x00 means the rest of the directory is empty
        if (entry->name[0] == 0x00) break;
        
        // 0xE5 means the file was deleted, skip it
        if ((unsigned char)entry->name[0] == 0xE5) continue;
        
        // Skip hidden files, system files, and long-filename entries (0x0F)
        if (entry->attributes == 0x0F || (entry->attributes & 0x02) || (entry->attributes & 0x04)) continue;
        
        // Print the short filename!
        vga_print("Found File: ");
        for(int j=0; j<8; j++) if(entry->name[j] != ' ') vga_print_char(entry->name[j]);
        vga_print_char('.');
        for(int j=0; j<3; j++) if(entry->ext[j] != ' ') vga_print_char(entry->ext[j]);
        vga_print("\n");
    }
}
uint32_t fat32_get_next_cluster(uint32_t current_cluster) {
    // 1 Sector = 512 Bytes. Each FAT entry is 4 Bytes (32 bits). 
    // Therefore, there are exactly 128 FAT entries per sector!
    uint32_t fat_sector = fat_start_sector + (current_cluster / 128);
    uint32_t fat_offset = (current_cluster % 128) * 4;
    
    uint8_t sector_buffer[512];
    ata_read_sector(fat_sector, sector_buffer);
    
    // Read the 32-bit FAT entry
    uint32_t next_cluster = *((uint32_t*)&sector_buffer[fat_offset]);
    
    // Mask out the top 4 bits (since it's FAT32, only 28 bits are used)
    return next_cluster & 0x0FFFFFFF; 
}

bool fat32_read_file(const char* target_name, const char* target_ext, uint8_t* out_buffer) {
    uint8_t root_buffer[512];
    uint32_t root_lba = cluster_to_lba(root_dir_cluster);
    ata_read_sector(root_lba, root_buffer);
    
    FAT32_DirectoryEntry* entries = (FAT32_DirectoryEntry*)root_buffer;
    uint32_t file_cluster = 0;
    
    // 1. Search the Root Directory for the file
    for (int i = 0; i < 16; i++) {
        FAT32_DirectoryEntry* entry = &entries[i];
        if (entry->name[0] == 0x00) break;
        if ((unsigned char)entry->name[0] == 0xE5) continue;
        if (entry->attributes == 0x0F || (entry->attributes & 0x08)) continue;
        
        // We use strict 8.3 FAT32 string matching
        if (str_match(entry->name, target_name, 8) && str_match(entry->ext, target_ext, 3)) {
            // Found it! Combine the high and low 16-bit halves to get the 32-bit starting cluster
            file_cluster = ((uint32_t)entry->first_cluster_high << 16) | entry->first_cluster_low;
            break;
        }
    }
    
    if (file_cluster == 0) {
        vga_print("Error: File not found!\n");
        return false;
    }
    
    // 2. Read the file, cluster by cluster
    uint32_t current_cluster = file_cluster;
    int buffer_offset = 0;
    
    // A cluster value of 0x0FFFFFF8 or higher means "End of File"
    while (current_cluster < 0x0FFFFFF8) { 
        uint32_t lba = cluster_to_lba(current_cluster);
        
        // Read all sectors in this cluster (mtools formats use 1 sector per cluster for tiny disks)
        for (int i = 0; i < sectors_per_cluster; i++) {
            ata_read_sector(lba + i, &out_buffer[buffer_offset]);
            buffer_offset += 512;
        }
        
        // Find the next cluster from the FAT table
        current_cluster = fat32_get_next_cluster(current_cluster);
    }
    
    return true;
}
int fat32_get_root_entries(FAT32_DirectoryEntry* out_buffer, int max_entries) {
    uint8_t root_buffer[512];
    uint32_t root_lba = cluster_to_lba(root_dir_cluster);
    ata_read_sector(root_lba, root_buffer);
    
    FAT32_DirectoryEntry* entries = (FAT32_DirectoryEntry*)root_buffer;
    int found_count = 0;
    
    for (int i = 0; i < 16; i++) {
        FAT32_DirectoryEntry* entry = &entries[i];
        if (entry->name[0] == 0x00) break;
        if ((unsigned char)entry->name[0] == 0xE5) continue;
        if (entry->attributes == 0x0F || (entry->attributes & 0x08)) continue;
        
        if (found_count < max_entries) {
            out_buffer[found_count] = *entry;
            found_count++;
        }
    }
    
    return found_count;
}

\end{lstlisting}

\subsection{fat32.h}
\begin{lstlisting}[language={C++}]
#pragma once
#include <stdint.h>

// Packed ensures the compiler doesn't add padding between variables!
// #pragma pack(push, 1)

// The FAT32 Volume Boot Record (VBR) located at Sector 0
struct FAT32_VBR {
    uint8_t  jump[3];
    char     oem_name[8];
    uint16_t bytes_per_sector;      // Usually 512
    uint8_t  sectors_per_cluster;
    uint16_t reserved_sectors;      // Boot record sectors before the FAT
    uint8_t  fat_count;             // Number of FATs (usually 2)
    uint16_t dir_entries;           // Unused in FAT32 (0)
    uint16_t total_sectors_16;      // Unused in FAT32 (0)
    uint8_t  media_descriptor;
    uint16_t sectors_per_fat_16;    // Unused in FAT32 (0)
    uint16_t sectors_per_track;
    uint16_t heads;
    uint32_t hidden_sectors;
    uint32_t total_sectors_32;      // Total sectors in volume
    
    // FAT32 Extended Boot Record
    uint32_t sectors_per_fat_32;    // Size of the FAT in sectors
    uint16_t flags;
    uint16_t fat_version;
    uint32_t root_cluster;          // Usually cluster 2
    uint16_t fsinfo_cluster;
    uint16_t backup_boot_sector;
    uint8_t  reserved[12];
    uint8_t  drive_number;
    uint8_t  nt_flags;
    uint8_t  signature;             // 0x28 or 0x29
    uint32_t serial_number;
    char     label[11];
    char     sys_identifier[8];     // "FAT32   "
    uint8_t  boot_code[420];
    uint16_t boot_signature;        // 0xAA55
}__attribute__((packed));

// A 32-byte FAT32 Directory Entry
struct FAT32_DirectoryEntry {
    char     name[8];               // Short filename (e.g., "HELLO   ")
    char     ext[3];                // Extension (e.g., "TXT")
    uint8_t  attributes;            // File, Directory, Hidden, Read-Only, etc.
    uint8_t  reserved;
    uint8_t  creation_time_tenths;
    uint16_t creation_time;
    uint16_t creation_date;
    uint16_t access_date;
    uint16_t first_cluster_high;    // High 16 bits of the data cluster
    uint16_t modify_time;
    uint16_t modify_date;
    uint16_t first_cluster_low;     // Low 16 bits of the data cluster
    uint32_t size;                  // File size in bytes
}__attribute__((packed));

// #pragma pack(pop)

void fat32_init();
void fat32_list_root();
bool fat32_read_file(const char* target_name, const char* target_ext, uint8_t* out_buffer);
int fat32_get_root_entries(FAT32_DirectoryEntry* out_buffer, int max_entries);



\end{lstlisting}

\subsection{gdt.cpp}
\begin{lstlisting}[language={C++}]
#include "gdt.h"
#include "tss.h"
extern struct tss_entry_t tss;
GDT gdt[6];
GDTR gdtr;

void set_entry(int index, uint32_t base, uint32_t limit, uint8_t access, uint8_t granularity){
    gdt[index].base_low = base & 0xFFFF;
    gdt[index].base_middle = (base >> 16) & 0xFF;
    gdt[index].base_high = (base >> 24) & 0xFF;
    gdt[index].limit_low = limit & 0xFFFF;
    gdt[index].access = access;
    gdt[index].granularity = ((limit >> 16) & 0x0F) | granularity;
}

extern "C" void gdt_flush(GDTR*);

void gdt_init(){
    set_entry(0,0,0,0,0);
    set_entry(1, 0, 0xFFFFFFFF, 0x9A, 0xCF);
    set_entry(2, 0, 0xFFFFFFFF, 0x92, 0xCF);
    set_entry(3, 0, 0xFFFFFFFF, 0xFA, 0xCF);
    set_entry(4, 0, 0xFFFFFFFF, 0xF2, 0xCF);
    uint32_t tss_base = (uint32_t)&tss;
    uint32_t tss_limit = sizeof(tss);
    set_entry(5, tss_base, tss_limit, 0xE9, 0x00);
    gdtr.limit = sizeof(gdt) - 1;
    gdtr.base = (uint32_t)&gdt;
    gdt_flush(&gdtr);
}


\end{lstlisting}

\subsection{gdt.h}
\begin{lstlisting}[language={C++}]
#pragma once
#pragma pack(1)
#include <stdint.h>
struct GDT{
    uint16_t limit_low, base_low;
    uint8_t base_middle, access, granularity, base_high;
};

struct GDTR{
    uint16_t limit;
    uint32_t base;
};

void gdt_init();
\end{lstlisting}

\subsection{idt.cpp}
\begin{lstlisting}[language={C++}]
#include "idt.h"


IDT idt[256];
IDTR idtr;

void idt_set_gate(int num, uint32_t base, uint16_t sel, uint8_t flag){
    idt[num].base_low = base & 0xFFFF;
    idt[num].base_high = (base >> 16) & 0xFFFF;
    idt[num].sel = sel;
    idt[num].zero = 0;
    idt[num].flags = flag;
}
extern "C" void isr0(); extern "C" void isr1(); extern "C" void isr2(); extern "C" void isr3();
extern "C" void isr4(); extern "C" void isr5(); extern "C" void isr6(); extern "C" void isr7();
extern "C" void isr8(); extern "C" void isr9(); extern "C" void isr10(); extern "C" void isr11();
extern "C" void isr12(); extern "C" void isr13(); extern "C" void isr14(); extern "C" void isr15();
extern "C" void isr16(); extern "C" void isr17(); extern "C" void isr18(); extern "C" void isr19();
extern "C" void isr20(); extern "C" void isr21(); extern "C" void isr22(); extern "C" void isr23();
extern "C" void isr24(); extern "C" void isr25(); extern "C" void isr26(); extern "C" void isr27();
extern "C" void isr28(); extern "C" void isr29(); extern "C" void isr30(); extern "C" void isr31();
extern "C" void isr128();

extern "C" void isr3();
extern "C" void isr32();
extern "C" void isr33();
extern "C" void isr44();
extern "C" void idt_flush(IDTR*);

void idt_init(){
    idtr.limit = sizeof(idt) - 1;
    idtr.base = (uint32_t)&idt;
    for(int i = 0; i<256; i++){
        idt_set_gate(i, 0, 0, 0);
    }
    idt_set_gate(0, (uint32_t)isr0, 0x08, 0x8E); idt_set_gate(1, (uint32_t)isr1, 0x08, 0x8E);
    idt_set_gate(2, (uint32_t)isr2, 0x08, 0x8E); idt_set_gate(3, (uint32_t)isr3, 0x08, 0x8E);
    idt_set_gate(4, (uint32_t)isr4, 0x08, 0x8E); idt_set_gate(5, (uint32_t)isr5, 0x08, 0x8E);
    idt_set_gate(6, (uint32_t)isr6, 0x08, 0x8E); idt_set_gate(7, (uint32_t)isr7, 0x08, 0x8E);
    idt_set_gate(8, (uint32_t)isr8, 0x08, 0x8E); idt_set_gate(9, (uint32_t)isr9, 0x08, 0x8E);
    idt_set_gate(10, (uint32_t)isr10, 0x08, 0x8E); idt_set_gate(11, (uint32_t)isr11, 0x08, 0x8E);
    idt_set_gate(12, (uint32_t)isr12, 0x08, 0x8E); idt_set_gate(13, (uint32_t)isr13, 0x08, 0x8E);
    idt_set_gate(14, (uint32_t)isr14, 0x08, 0x8E); idt_set_gate(15, (uint32_t)isr15, 0x08, 0x8E);
    idt_set_gate(16, (uint32_t)isr16, 0x08, 0x8E); idt_set_gate(17, (uint32_t)isr17, 0x08, 0x8E);
    idt_set_gate(18, (uint32_t)isr18, 0x08, 0x8E); idt_set_gate(19, (uint32_t)isr19, 0x08, 0x8E);
    idt_set_gate(20, (uint32_t)isr20, 0x08, 0x8E); idt_set_gate(21, (uint32_t)isr21, 0x08, 0x8E);
    idt_set_gate(22, (uint32_t)isr22, 0x08, 0x8E); idt_set_gate(23, (uint32_t)isr23, 0x08, 0x8E);
    idt_set_gate(24, (uint32_t)isr24, 0x08, 0x8E); idt_set_gate(25, (uint32_t)isr25, 0x08, 0x8E);
    idt_set_gate(26, (uint32_t)isr26, 0x08, 0x8E); idt_set_gate(27, (uint32_t)isr27, 0x08, 0x8E);
    idt_set_gate(28, (uint32_t)isr28, 0x08, 0x8E); idt_set_gate(29, (uint32_t)isr29, 0x08, 0x8E);
    idt_set_gate(30, (uint32_t)isr30, 0x08, 0x8E); idt_set_gate(31, (uint32_t)isr31, 0x08, 0x8E);
    idt_set_gate(3, (uint32_t)isr3, 0x08, 0x8E);
    idt_set_gate(32, (uint32_t)isr32, 0x08, 0x8E);
    idt_set_gate(33, (uint32_t)isr33, 0x08, 0x8E);
    idt_set_gate(128, (uint32_t)isr128, 0x08, 0xEE);
    idt_set_gate(44, (uint32_t)isr44, 0x08, 0x8E);
    idt_flush(&idtr);
}
\end{lstlisting}

\subsection{idt.h}
\begin{lstlisting}[language={C++}]
#pragma once
#pragma pack(1)
#include <stdint.h>

struct IDT{
    uint16_t base_low, sel;
    uint8_t zero, flags;
    uint16_t base_high;
};

struct IDTR{
    uint16_t limit;
    uint32_t base;
};

void idt_init();
\end{lstlisting}

\subsection{initrd.cpp}
\begin{lstlisting}[language={C++}]
#include "initrd.h"
#include "vfs.h"
#include "string.h"
#include "vga.h"


static uint32_t octal_to_int(const char* s) {
    uint32_t result = 0;
    while (*s) {
        result = result * 8 + (*s - '0');
        s++;
    }
    return result;
}

struct tar_header{
    char name[100];
    char mode[8];
    char uid[8];
    char gid[8];
    char size[12];
    char mtime[12];
    char checksum[8];
    char type;
    char linkname[100];
    char padding[255];
}__attribute__((packed));

static fs_node nodes[64];
static tar_header* headers[64];
static uint32_t file_count = 0;

static uint32_t tarfs_read(fs_node* node, uint32_t offset, uint32_t size, char* buf) {
    tar_header* hdr = headers[node->inode];
    char* content = (char*)hdr + 512;
    
    if (offset > node->length) return 0;
    if (offset + size > node->length) size = node->length - offset;
    
    for (uint32_t i = 0; i < size; i++) {
        buf[i] = content[offset + i];
    }
    return size;
}


fs_node* initrd_init(uint32_t mod_start){
    tar_header* hdr = (tar_header*) mod_start;
    while(hdr->name[0] != '\0'){
        uint32_t file_size = octal_to_int(hdr->size);
        for(int i = 0; i<100; i++){
            nodes[file_count].name[i] = hdr->name[i];
        }
        nodes[file_count].flags = VFS_FILE;
        nodes[file_count].length = file_size;
        nodes[file_count].inode = file_count;
        nodes[file_count].read = tarfs_read;
        headers[file_count] = hdr;
        file_count++;
        uint32_t blocks = (file_size + 511) / 512;
        hdr = (tar_header*)((char*)hdr + 512 + blocks * 512);
    }
    return nodes;
}
fs_node* initrd_findfile(const char* file){
    for(uint32_t i = 0; i<file_count; i++){
        if(strcmp(nodes[i].name, file)){
            return &nodes[i];
        }
    }
    return 0;
}
void initrd_listfiles(){
    for(uint32_t i = 0; i<file_count; i++){
        vga_print(nodes[i].name);
        vga_print(" ");
    }
    vga_print("\n");
}
\end{lstlisting}

\subsection{initrd.h}
\begin{lstlisting}[language={C++}]
#pragma once
#include <stdint.h>
#include "vfs.h"

fs_node* initrd_init(uint32_t mod_start);
fs_node* initrd_findfile(const char* name);
void initrd_listfiles();
\end{lstlisting}

\subsection{io.h}
\begin{lstlisting}[language={C++}]
#pragma once
#include <stdint.h>

static inline void outb(uint16_t port, uint8_t val){
    asm volatile ( "outb %0, %1" : : "a"(val), "Nd"(port) );
}

static inline uint8_t inb(uint16_t port) {
    uint8_t ret;
    asm volatile ( "inb %1, %0" : "=a"(ret) : "Nd"(port) );
    return ret;
}

static inline void io_wait(void) {
    outb(0x80, 0);
}

static inline void insw(uint16_t port, void* addr, uint32_t count) {
    asm volatile("cld; rep insw" : "+D"(addr), "+c"(count) : "d"(port) : "memory");
}

static inline void outsw(uint16_t port, const void* addr, uint32_t count) {
    asm volatile("cld; rep outsw" : "+S"(addr), "+c"(count) : "d"(port) : "memory");
}

\end{lstlisting}

\subsection{isr.cpp}
\begin{lstlisting}[language={C++}]
#include "vga.h"
#include "pic.h"
#include "io.h"
#include "keyboard.h"
#include "string.h"
#include "memory.h"
#include "timer.h"
#include "initrd.h"

struct regs_t {
    uint32_t edi, esi, ebp, esp, ebx, edx, ecx, eax;
    uint32_t int_no, err_code;                 
    uint32_t eip, cs, eflags, useresp, ss;  
};

// const char* message = "The Abyss is staring at you, mortal!!\n";
extern char* message;
const char* exception_messages[32] = {
    "Division By Zero", "Debug", "Non Maskable Interrupt", "Breakpoint",
    "Into Detected Overflow", "Out of Bounds", "Invalid Opcode", "No Coprocessor",
    "Double Fault", "Coprocessor Segment Overrun", "Bad TSS", "Segment Not Present",
    "Stack Fault", "General Protection Fault", "Page Fault", "Unknown Interrupt",
    "Coprocessor Fault", "Alignment Check", "Machine Check", "Reserved", 
    "Reserved", "Reserved", "Reserved", "Reserved", "Reserved", 
    "Reserved", "Reserved", "Reserved", "Reserved", "Reserved", 
    "Reserved", "Reserved"
};

extern "C" void fault_handler(const struct regs_t* r){
    if(r->int_no < 32){
        vga_clear();
        vga_print("Abyss is angered!! Your connection with Abyss is severed!!\n");
        vga_print(exception_messages[r->int_no]);
    }else if(r->int_no == 128){
        if(r->eax == 1){
            char* add = (char*)r->ebx;
            vga_print(add);
        }

    }
    while(1){
        asm volatile("cli");
        asm volatile("hlt");
    }
}
extern "C" void timer_tick(); 
char keyboard_buffer[256];
int buffer_ind = 0;
void execute_command(){
    keyboard_buffer[buffer_ind] = '\0';
    vga_print("\n");
    if(strcmp(keyboard_buffer, "clear")){
        vga_clear();
        vga_print(message);
    }
    else if(strcmp(keyboard_buffer, "hail")){
        vga_print("THE OLDEST DARKNESS\n");
    }
    else if(strcmp(keyboard_buffer, "help")){
        vga_print("These are the helping commands mortal :\n 1. clear\n 2. hail\n 3. alloc/malloc\n 4. uptime\n 5. cat <filename>\n 6. help\n 7. ls\n");
    }
    else if(strcmp(keyboard_buffer, "alloc") or strcmp(keyboard_buffer, "malloc")){
        char* dynamic_string = (char*)malloc(20);
        dynamic_string[0] = 'C';
        dynamic_string[1] = 'o';
        dynamic_string[2] = 'r';
        dynamic_string[3] = 'r';
        dynamic_string[4] = 'o';
        dynamic_string[5] = 'd';
        dynamic_string[6] = 'e';
        dynamic_string[7] = '\0';
        vga_print("The Abyss gave you some space and stated:\n");
        vga_print(dynamic_string);
        vga_print("\n");
    }
    else if(strcmp(keyboard_buffer, "uptime")){
        uint32_t time = timer_get_ticks();
        vga_print("The Abyss has been rampant for ");
        char ch[256];
        itoa(time, ch);
        vga_print(ch);
        vga_print(" ticks!!\n");
    }else if(keyboard_buffer[0] == 'c' and keyboard_buffer[1] == 'a' and keyboard_buffer[2] == 't' and keyboard_buffer[3] == ' '){
        char* filename = keyboard_buffer + 4;
        fs_node* file = initrd_findfile(filename);
        if(file){
            char buf[512] = {0};
            vfs_read(file, 0, file->length, buf);
            vga_print(buf);
            // vga_print("\n");
        }else{
            vga_print("No file named ");
            vga_print(filename);
            vga_print(" in inventory, mortal!!\n");
        }
    }else if(strcmp(keyboard_buffer, "ls")){
        initrd_listfiles();
    }
    else if(buffer_ind > 0){
        vga_print("Abyss knows no such thing as ");
        vga_print("' ");
        vga_print(keyboard_buffer);
        vga_print(" '");
        vga_print("\n");
    }
    buffer_ind = 0;
    vga_print("Abyss >> ");
}
extern "C" void isr3_handler(){
    vga_print("Interrupt 3 fired!!");
}
extern "C" void isr32_handler() {
    timer_tick();
    pic_send_eoi(0);
}

extern "C" void isr33_handler() {
    uint8_t scancode = inb(0x60);
    if (!(scancode & 0x80)) { 
        char c = kbdus[scancode]; 
        
        if (c == '\n') {
            execute_command(); 
        } 
        else if (c == '\b') {
            if (buffer_ind> 0) {
                buffer_ind--;
                vga_print_char('\b');
            }
        } 
        else if (c != 0 && buffer_ind < 255) { 
            keyboard_buffer[buffer_ind] = c;
            buffer_ind++;
            vga_print_char(c);
        }
    }
    pic_send_eoi(1); 
}

\end{lstlisting}

\subsection{kernel.cpp}
\begin{lstlisting}[language={C++}]
#include "vga.h"
#include "gdt.h"
#include "idt.h"
#include "pic.h"
#include "paging.h"
#include "tss.h"
#include "memory.h"
#include "timer.h"
#include "task.h"
#include "multiboot.h"
#include "vfs.h"
#include "initrd.h"
#include "vesa.h"
#include "mouse.h"
#include "ui.h"
#include "ata.h"
#include "fat32.h"

extern "C" void jump_usermode();
extern "C" void test_user_function(){
    const char* str = "Hello from User Mode via System Call!\n";
    asm volatile(
        "mov %0, %%eax;"
        "mov %1, %%ebx;"
        "int $0x80;"
        : : "r"(1), "r"(str) : "eax", "ebx"
    );
    
    while(1){}
}
const char* message = "The Abyss is staring at you, mortal!!\n";
extern "C" void perform_task_switch();
void task_a(){
    while(1){
        vga_print("A ");
    }
}
void task_b(){
    while(1){
        vga_print("B ");
    }
}
extern "C" void kernel_main(multiboot_info* mbd){
    gdt_init();
    uint32_t esp;
    asm volatile("mov %%esp, %0" : "=r"(esp));
    tss_init(esp);
    idt_init();
    pic_remap(32, 40);
    init_paging();
    memory_init();
    timer_init(100);
    init_tasking();
    vesa_init((uint32_t)mbd->framebuffer_addr, mbd->framebuffer_width, mbd->framebuffer_height, mbd->framebuffer_pitch);
    mouse_init();
    ata_init();
    uint8_t disk_buf[512];
    for(int i = 0; i<512; i++) disk_buf[i] = 0;
    const char* disk_msg = "The Hard drive is ready mortal!";
    for(int i = 0; disk_msg[i] != '\0'; i++){
        disk_buf[i] = disk_msg[i];
    }
    ata_write_sector(1, disk_buf);
    vga_clear();
    draw_rect(0, 0, 1024, 768, 0x000000);
    fat32_init();
    // fat32_list_root();

    // uint8_t file_buffer[2048];
    // for(int i = 0; i<2048; i++) file_buffer[i] = 0;
    // if (fat32_read_file("HELLO   ", "TXT", file_buffer)) {
    //     vga_print("\nRead from physical Hard Drive:\n> ");
    //     vga_print((const char*)file_buffer);
    //     vga_print("\n");
    // }

    // vesa_swap_buffers();
    while(1){
        ui_render();
        vesa_swap_buffers();
    }
    // draw_gradient_rect(0,0,1024,768, 0x100010, 0x004554);
    vga_print("Abyss v1.0 loading...\nLoading Virtual File System...\nFiles loaded...\n...\n...\nHello Mortal...");
    if (mbd->mods_count > 0) {
        multiboot_mod_list* mod = (multiboot_mod_list*)mbd->mods_addr;
        initrd_init(mod->mod_start);
        const char* name = "hell.txt";
        fs_node* file = initrd_findfile(name);
        if (file) {
            char buf[128] = {0};
            vfs_read(file, 0, file->length, buf);
            vga_print("VFS Read: ");
            vga_print(buf);
            vga_print("\n");
        }else{
            vga_print("Filename ");
            vga_print(name);
            vga_print(" not found!!\n");
        }
    }else{
        vga_print("ERROR : no initrd found!!");
    }
    vga_print(message);
    vga_print("You want help? HAHA!! Type 'help', mortal\n");
    vga_print("Abyss >> ");
    // Task* t1 = (Task*)malloc(sizeof(Task));
    // Task* t2 = (Task*)malloc(sizeof(Task));

    // create_task(t1, task_a);
    // create_task(t2, task_b);
    asm volatile ("sti");
    while(1){
        ui_render();
        asm volatile ("hlt");
    }
    // while(1){}
    // asm volatile ("int $0x8");
    // asm volatile ("sti");
    // jump_usermode();
    // while(1){}
}
\end{lstlisting}

\subsection{keyboard.h}
\begin{lstlisting}[language={C++}]
#pragma once

const char kbdus[128] = {
    0,  27, '1', '2', '3', '4', '5', '6', '7', '8', /* 9 */
  '9', '0', '-', '=', '\b', /* Backspace */
  '\t', /* Tab */
  'q', 'w', 'e', 'r', /* 19 */
  't', 'y', 'u', 'i', 'o', 'p', '[', ']', '\n', /* Enter key */
    0,  /* 29   - Control */
  'a', 's', 'd', 'f', 'g', 'h', 'j', 'k', 'l', ';', /* 39 */
 '\'', '`',   0,  /* Left shift */
 '\\', 'z', 'x', 'c', 'v', 'b', 'n', /* 49 */
  'm', ',', '.', '/',   0,  /* Right shift */
  '*',
    0,  /* Alt */
  ' ',  /* Space bar */
    0,  /* Caps lock */
    0, 0, 0, 0, 0, 0, 0, 0, 0, 0, /* F1 to F10 keys */
    0,  /* Num lock */
    0,  /* Scroll Lock */
    0,  /* Home key */
    0,  /* Up Arrow */
    0,  /* Page Up */
  '-',
    0,  /* Left Arrow */
    0,
    0,  /* Right Arrow */
  '+',
    0,  /* End key*/
    0,  /* Down Arrow */
    0,  /* Page Down */
    0,  /* Insert Key */
    0,  /* Delete Key */
    0, 0, 0,
    0,  /* F11 Key */
    0,  /* F12 Key */
    0, /* All other keys are undefined */
};

\end{lstlisting}

\subsection{memory.cpp}
\begin{lstlisting}[language={C++}]
#include "memory.h"
#include "vga.h"

uint32_t heap_start = 0x200000;
uint32_t current_heap_pointer;

void memory_init(){
    current_heap_pointer = heap_start;
}
void* malloc(size_t size){
    uint32_t allocated_address = current_heap_pointer;
    current_heap_pointer+=size;
    if(current_heap_pointer >= 0x2000000){
        vga_print("The Abyss is filled with your trash, mortal\n");
        while(1){
            asm volatile("cli; hlt");
        }
    }
    return (void*)allocated_address;
}
\end{lstlisting}

\subsection{memory.h}
\begin{lstlisting}[language={C++}]
#pragma once
#include <stdint.h>
#include <stddef.h>

void memory_init();
void* malloc(size_t size);
\end{lstlisting}

\subsection{mouse.cpp}
\begin{lstlisting}[language={C++}]
#include "mouse.h"
#include "io.h"
#include "vesa.h"
#include "pic.h"
#include "ui.h" // We will create this next!

int mouse_x = 512;
int mouse_y = 384;
uint8_t mouse_cycle = 0;
uint8_t mouse_byte[3];
bool mouse_left_pressed = false; 

const int CURSOR_W = 10;
const int CURSOR_H = 15;
const uint8_t cursor_bitmap[15][10] = {
    {2,0,0,0,0,0,0,0,0,0},
    {2,2,0,0,0,0,0,0,0,0},
    {2,1,2,0,0,0,0,0,0,0},
    {2,1,1,2,0,0,0,0,0,0},
    {2,1,1,1,2,0,0,0,0,0},
    {2,1,1,1,1,2,0,0,0,0},
    {2,1,1,1,1,1,2,0,0,0},
    {2,1,1,1,1,1,1,2,0,0},
    {2,1,1,1,1,1,1,1,2,0},
    {2,1,1,1,1,1,1,1,1,2},
    {2,1,1,1,1,2,2,2,2,2},
    {2,1,1,2,1,2,0,0,0,0},
    {2,1,2,0,2,1,2,0,0,0},
    {2,2,0,0,2,1,2,0,0,0},
    {2,0,0,0,0,2,2,0,0,0}
};

void mouse_wait(uint8_t type){
    uint32_t timeout = 100000;
    if (type == 0) {
        while (timeout--) { if ((inb(0x64) & 1) == 1) return; }
    } else {
        while (timeout--) { if ((inb(0x64) & 2) == 0) return; }
    }
}

void mouse_write(uint8_t a_write) {
    mouse_wait(1);
    outb(0x64, 0xD4); 
    mouse_wait(1);
    outb(0x60, a_write);
}

uint8_t mouse_read() {
    mouse_wait(0);
    return inb(0x60);
}

void mouse_init() {
    uint8_t status;
    mouse_wait(1);
    outb(0x64, 0xA8);
    mouse_wait(1);
    outb(0x64, 0x20); 
    mouse_wait(0);
    status = (inb(0x60) | 2);
    mouse_wait(1);
    outb(0x64, 0x60);
    mouse_wait(1);
    outb(0x60, status);
    mouse_write(0xF6);
    mouse_read();
    mouse_write(0xF4);
    mouse_read();
}

void mouse_draw() {
    for (int i = 0; i < CURSOR_H; i++) {
        for (int j = 0; j < CURSOR_W; j++) {
            if (cursor_bitmap[i][j] == 1) {
                put_pixel(mouse_x + j, mouse_y + i, 0x000000); 
            } else if (cursor_bitmap[i][j] == 2) {
                put_pixel(mouse_x + j, mouse_y + i, 0xFFFFFF); 
            }
        }
    }
}

extern "C" void isr44_handler() {
    uint8_t status = inb(0x64);
    if ((status & 0x20) == 0) {
        pic_send_eoi(12);
        return; 
    }
    
    mouse_byte[mouse_cycle] = inb(0x60); 
    if (mouse_cycle == 0 && (mouse_byte[0] & 0x08) == 0) {
        pic_send_eoi(12);
        return;
    }
    
    mouse_cycle++; 
    if (mouse_cycle == 3) {
        mouse_cycle = 0; 
        mouse_left_pressed = (mouse_byte[0] & 0x01);
        
        int d_x = (int8_t)mouse_byte[1]; 
        int d_y = (int8_t)mouse_byte[2]; 
        
        mouse_x += d_x;
        mouse_y -= d_y; 
        
        if (mouse_x < 0) mouse_x = 0;
        if (mouse_y < 0) mouse_y = 0;
        if (mouse_x > 1023 - CURSOR_W) mouse_x = 1023 - CURSOR_W; 
        if (mouse_y > 767 - CURSOR_H) mouse_y = 767 - CURSOR_H;   
        
        // Signal the UI to draw a new frame!
        needs_redraw = true;
    }
    pic_send_eoi(12); 
}

\end{lstlisting}

\subsection{mouse.h}
\begin{lstlisting}[language={C++}]
#pragma once
#include <stdint.h>

extern int mouse_x;
extern int mouse_y;
extern uint8_t mouse_byte[3];
extern bool mouse_left_pressed;

void mouse_init();
void mouse_draw();

\end{lstlisting}

\subsection{multiboot.h}
\begin{lstlisting}[language={C++}]
#pragma once
#include <stdint.h>

struct multiboot_info {
    uint32_t flags;
    uint32_t mem_lower;
    uint32_t mem_upper;
    uint32_t boot_device;
    uint32_t cmdline;
    uint32_t mods_count;
    uint32_t mods_addr;
    uint32_t num;
    uint32_t size;
    uint32_t addr;
    uint32_t shndx;
    uint32_t mmap_length;
    uint32_t mmap_addr;
    uint32_t drives_length;
    uint32_t drives_addr;
    uint32_t config_table;
    uint32_t boot_loader_name;
    uint32_t apm_table;

    uint32_t vbe_control_info;
    uint32_t vbe_mode_info;
    uint16_t vbe_mode;
    uint16_t vbe_interface_seg;
    uint16_t vbe_interface_off;
    uint16_t vbe_interface_len;
    
    uint64_t framebuffer_addr;    // The physical RAM address where the screen starts
    uint32_t framebuffer_pitch;   // How many BYTES make up one horizontal row of pixels
    uint32_t framebuffer_width;   // Screen width in pixels (1024)
    uint32_t framebuffer_height;  // Screen height in pixels (768)
    uint8_t  framebuffer_bpp;     // Bits Per Pixel (32 = 4 bytes per pixel)
    uint8_t  framebuffer_type;    // 1 = RGB color mode
    uint8_t  color_info[6];       // Which bits are Red, Green, Blue
} __attribute__((packed));

struct multiboot_mod_list {
    uint32_t mod_start; 
    uint32_t mod_end; 
    uint32_t cmdline;
    uint32_t pad;
} __attribute__((packed));

\end{lstlisting}

\subsection{paging.cpp}
\begin{lstlisting}[language={C++}]
#include "paging.h"

struct page_directory_t page_directory __attribute__((aligned(4096)));
struct page_table_t first_page_table __attribute__((aligned(4096)));

extern "C" void load_page_directory(uint32_t*);
extern "C" void enable_paging();

// void init_paging(){
//     for(int i = 0; i<1024; i++){
//         page_directory.tablesPhysical[i] = 0x02;
//     }
//     for(int i = 0; i<1024; i++){
//         first_page_table.page[i] = (i * 4096) | 7;
//     }
//     page_directory.tablesPhysical[0] = ((uint32_t)&first_page_table) | 7;
//     load_page_directory(page_directory.tablesPhysical);
//     enable_paging();
// }
void init_paging() {

    for (int i = 0; i < 1024; i++) {
        page_directory.tablesPhysical[i] = (i << 22) | 0x83;
    }
    load_page_directory(page_directory.tablesPhysical);
    enable_paging();
}

\end{lstlisting}

\subsection{paging.h}
\begin{lstlisting}[language={C++}]
#pragma once
#include <stdint.h>

struct page_table_t{
    uint32_t page[1024];
};
struct page_directory_t{
    uint32_t tablesPhysical[1024];
    struct page_table_t* tables[1024];
};
void init_paging();
\end{lstlisting}

\subsection{pic.cpp}
\begin{lstlisting}[language={C++}]
#include "pic.h"
#include "io.h"

#define PIC1_COMMAND 0x20
#define PIC1_DATA    0x21
#define PIC2_COMMAND 0xA0
#define PIC2_DATA    0xA1

#define ICW1_INIT    0x10
#define ICW1_ICW4    0x01
#define ICW4_8086    0x01

void pic_remap(int offset1, int offset2) {
    unsigned char a1, a2;
    
    // Save current masks
    a1 = inb(PIC1_DATA);
    a2 = inb(PIC2_DATA);
    
    // Start initialization sequence in cascade mode
    outb(PIC1_COMMAND, ICW1_INIT | ICW1_ICW4);
    io_wait();
    outb(PIC2_COMMAND, ICW1_INIT | ICW1_ICW4);
    io_wait();
    
    // Set offsets (this is where we map it to 32 and 40)
    outb(PIC1_DATA, offset1);
    io_wait();
    outb(PIC2_DATA, offset2);
    io_wait();
    
    // Tell Master PIC that there is a Slave PIC at IRQ2 (0000 0100)
    outb(PIC1_DATA, 4);
    io_wait();
    
    // Tell Slave PIC its cascade identity (0000 0010)
    outb(PIC2_DATA, 2);
    io_wait();
    
    // 8086/88 (MCS-80/85) mode
    outb(PIC1_DATA, ICW4_8086);
    io_wait();
    outb(PIC2_DATA, ICW4_8086);
    io_wait();
    
    // Restore saved masks
    outb(PIC1_DATA, a1);
    outb(PIC2_DATA, a2);
}

// End of Interrupt (EOI) command
void pic_send_eoi(unsigned char irq) {
    if(irq >= 8)
        outb(PIC2_COMMAND, 0x20); // Send to Slave
    
    outb(PIC1_COMMAND, 0x20); // Send to Master
}

\end{lstlisting}

\subsection{pic.h}
\begin{lstlisting}[language={C++}]
#pragma once
#include <stdint.h>

void pic_remap(int offset1, int offset2);
void pic_send_eoi(unsigned char irq);
\end{lstlisting}

\subsection{string.cpp}
\begin{lstlisting}[language={C++}]
#include "string.h"

bool strcmp(const char* a, const char* b){
    int i = 0;
    while(a[i] != '\0' and b[i] != '\0'){
        if(a[i] != b[i]){
            return false;
        }
        i++;
    }
    return (a[i] == '\0' and b[i] == '\0');
}
void itoa(int n, char st[]){
    int i = 0;
    bool is_neg = 0;
    if(n == 0){
        st[i++] = '0';
        st[i] = '\0';
        return;
    }
    if(n < 0){
        n = -n;
        is_neg = 1;
    }
    while(n != 0){
        int r = n % 10;
        n/=10;
        st[i++] = '0' + r;
    }
    if(is_neg){
        st[i++] = '-';
    }
    st[i] = '\0';
    int start = 0, end = i-1;
    while(start < end){
        char temp = st[start];
        st[start] = st[end];
        st[end] = temp;
        start++;
        end--;
    }
}
\end{lstlisting}

\subsection{string.h}
\begin{lstlisting}[language={C++}]
#pragma once
#include <stdint.h>

bool strcmp(const char* a, const  char* b);
void itoa(int n, char st[]);
\end{lstlisting}

\subsection{task.cpp}
\begin{lstlisting}[language={C++}]
#include "task.h"
#include "memory.h"
#include "vga.h"

Task* current_task;
Task* ready_queue;
uint32_t next_pid = 1;

void init_tasking(){
    asm volatile ("cli");
    current_task = (Task*)malloc(sizeof(Task));
    current_task->id = next_pid;
    next_pid++;
    current_task->esp = 0;
    current_task->ebp = 0;
    current_task->eip = 0;
    current_task->next = current_task;
    ready_queue = current_task;
    asm volatile ("sti");
}  

void create_task(Task* new_task, void (*main)()){
    uint32_t st = (uint32_t)malloc(4096);
    uint32_t* stack = (uint32_t*)(st + 4096);
    *(--stack) = 0x202;
    *(--stack) = 0x08;
    *(--stack) = (uint32_t)main;
    for(int i = 0; i<8; i++){
        *(--stack) = 0;
    }
    new_task->esp = (uint32_t)stack;
    new_task->ebp = (uint32_t)stack;
    new_task->eip = (uint32_t)main;
    new_task->id = next_pid++;
    new_task->next = current_task->next;
    current_task->next = new_task;
}
\end{lstlisting}

\subsection{task.h}
\begin{lstlisting}[language={C++}]
#pragma once
#include <stdint.h>

struct Task{
    uint32_t id;
    uint32_t esp, ebp;
    uint32_t eip;
    Task* next;
};

void init_tasking();
void create_task(Task* t, void (*main)());
void switch_task();
\end{lstlisting}

\subsection{timer.cpp}
\begin{lstlisting}[language={C++}]
#include "timer.h"
#include "io.h"

uint32_t tick_count = 0;
void timer_init(uint32_t frequency){
    uint32_t divisor = 1193180 / frequency;
    outb(0x43, 0x36);
    uint8_t low = (uint32_t)(divisor & 0xFF);
    uint8_t high = (uint32_t)((divisor >> 8) & 0xFF);
    outb(0x40, low);
    outb(0x40, high);
}
uint32_t timer_get_ticks(){
    return tick_count;
}
extern "C" void timer_tick(){
    tick_count++;
}
\end{lstlisting}

\subsection{timer.h}
\begin{lstlisting}[language={C++}]
#pragma once
#include <stdint.h>

void timer_init(uint32_t frequency);
uint32_t timer_get_ticks();
\end{lstlisting}

\subsection{tss.cpp}
\begin{lstlisting}[language={C++}]
#include "tss.h"

struct tss_entry_t tss;

void tss_init(uint32_t kernel_stack) {
    uint8_t* tss_bytes = (uint8_t*)&tss;
    for (uint32_t i = 0; i < sizeof(tss); i++) {
        tss_bytes[i] = 0;
    }
    tss.ss0 = 0x10;
    tss.esp0 = kernel_stack;
    tss.iomap_base = sizeof(tss);
    asm volatile("ltr %%ax" : : "a"(0x2B));
}

\end{lstlisting}

\subsection{tss.h}
\begin{lstlisting}[language={C++}]
#pragma once
#include <stdint.h>

struct tss_entry_t {
    uint32_t prev_tss;
    uint32_t esp0; // The ring 0 stack pointer!
    uint32_t ss0;  // The ring 0 stack segment!
    uint32_t esp1;
    uint32_t ss1;
    uint32_t esp2;
    uint32_t ss2;
    uint32_t cr3;
    uint32_t eip;
    uint32_t eflags;
    uint32_t eax, ecx, edx, ebx, esp, ebp, esi, edi;
    uint32_t es, cs, ss, ds, fs, gs;
    uint32_t ldt;
    uint16_t trap;
    uint16_t iomap_base;
} __attribute__((packed));

void tss_init(uint32_t kernel_stack);

\end{lstlisting}

\subsection{ui.cpp}
\begin{lstlisting}[language={C++}]
#include "ui.h"
#include "vesa.h"
#include "mouse.h"
#include "fat32.h"

bool needs_redraw = true; // Always draw the very first frame

Window windows[3] = {
    {100, 100, 400, 300, "Abyss OS Settings", false, 0, 0},
    {550, 250, 400, 200, "System Monitor", false, 0, 0},
    {200, 150, 350, 250, "File Explorer", false, 0, 0}
};
int num_windows = 3;

// int num_windows = 2;

void draw_window(uint32_t x, uint32_t y, uint32_t width, uint32_t height, const char* title) {
    // Window Body
    draw_rect(x, y, width, height, 0xCCCCCC); // Light gray
    
    // Window Borders (Creates a slight 90s 3D pop-out effect)
    draw_rect(x, y, width, 1, 0xFFFFFF); // Top
    draw_rect(x, y, 1, height, 0xFFFFFF); // Left
    draw_rect(x + width - 1, y, 1, height, 0x555555); // Right
    draw_rect(x, y + height - 1, width, 1, 0x555555); // Bottom
    
    // Title Bar
    draw_rect(x + 2, y + 2, width - 4, 20, 0x0000AA); // Dark blue title bar
    draw_string(title, x + 6, y + 8, 0xFFFFFF, 0x0000AA, 2); // White text
}

bool is_any_window_dragged = false;

void ui_update() {
    // 1. If the mouse is released, drop everything
    if (!mouse_left_pressed) {
        for (int i = 0; i < num_windows; i++) {
            windows[i].is_dragged = false;
        }
        is_any_window_dragged = false;
        return;
    }

    // 2. If we are currently dragging a window, ONLY update that specific window
    if (is_any_window_dragged) {
        for (int i = 0; i < num_windows; i++) {
            Window* w = &windows[i];
            if (w->is_dragged) {
                w->x = mouse_x - w->drag_offset_x;
                w->y = mouse_y - w->drag_offset_y;
                needs_redraw = true;
                return; // Exit early so we don't grab other windows!
            }
        }
    }

    // 3. If we aren't dragging anything yet, check for a NEW click!
    for (int i = num_windows - 1; i >= 0; i--) {
        Window* w = &windows[i];
        
        // Check collision with the ENTIRE WINDOW (so clicking the body brings it to front)
        if (mouse_x >= w->x && mouse_x <= w->x + w->width &&
            mouse_y >= w->y && mouse_y <= w->y + w->height) {
            
            // Bring to front! (If it's not already on top)
            if (i != num_windows - 1) {
                Window clicked_window = windows[i];
                
                // Shift all windows above it down by 1 position
                for (int j = i; j < num_windows - 1; j++) {
                    windows[j] = windows[j + 1];
                }
                
                // Put our clicked window at the very end of the array (the top)
                windows[num_windows - 1] = clicked_window;
                
                // Re-point 'w' to the new top window so dragging works
                w = &windows[num_windows - 1];
            }
            
            // Check if they specifically clicked the Title Bar to start dragging
            if (mouse_y <= w->y + 24) {
                w->is_dragged = true;
                w->drag_offset_x = mouse_x - w->x;
                w->drag_offset_y = mouse_y - w->y;
                is_any_window_dragged = true;
            }
            
            needs_redraw = true; // Force redraw to show it popped to the front!
            break; 
        }
    }
}



void ui_render() {
    ui_update(); // Process mouse clicks first!
    
    if (!needs_redraw) return;
    // 1. Draw Background
    draw_gradient_rect(0, 0, 1024, 768, 0x100010, 0x004554);
    
    // 2. Draw Windows from our array
    for (int i = 0; i < num_windows; i++) {
        Window* w = &windows[i];
        draw_window(w->x, w->y, w->width, w->height, w->title);
        if(w->title[0] == 'F' and w->title[1] == 'i'){
            draw_rect(w->x + 2, w->y + 24, w->width - 4, w->height - 26, 0xFFFFFF);
            FAT32_DirectoryEntry files[16];
            int file_count = fat32_get_root_entries(files, 16);
            int file_y = w->y + 30;
            for (int f = 0; f < file_count; f++) {
                draw_rect(w->x + 10, file_y, 16, 16, 0x0055AA);
                
                // Construct the filename string (8 chars name + '.' + 3 chars ext)
                char filename[13];
                int idx = 0;
                for(int j=0; j<8; j++) if(files[f].name[j] != ' ') filename[idx++] = files[f].name[j];
                filename[idx++] = '.';
                for(int j=0; j<3; j++) if(files[f].ext[j] != ' ') filename[idx++] = files[f].ext[j];
                filename[idx] = '\0';
                
                // Draw the text in black
                draw_string(filename, w->x + 35, file_y + 4, 0x000000, 0xFFFFFF, 1);
                
                file_y += 24;
            }
        }
    }
    // 3. Draw Mouse
    mouse_draw();
    // 4. Swap buffers
    vesa_swap_buffers();
    
    needs_redraw = false;
}



\end{lstlisting}

\subsection{ui.h}
\begin{lstlisting}[language={C++}]
#pragma once
#include <stdint.h>

struct Window {
    int x, y;
    int width, height;
    const char* title;
    
    // State for dragging
    bool is_dragged;
    int drag_offset_x;
    int drag_offset_y;
};

extern bool needs_redraw;
void ui_render();

\end{lstlisting}

\subsection{vesa.cpp}
\begin{lstlisting}[language={C++}]
#include "vesa.h"
#include "font8x8.h"
#include "memory.h"

// #include <stdlib.h>

uint32_t* framebuffer;
uint32_t* backbuffer;
static uint32_t  fb_width;
static uint32_t  fb_height;
static uint32_t  fb_pitch;

void vesa_init(uint32_t addr, uint32_t width, uint32_t height, uint32_t pitch) {
    framebuffer = (uint32_t*)addr;
    fb_width    = width;
    fb_height   = height;
    fb_pitch    = pitch;

    backbuffer = (uint32_t*)malloc(width * height * 4);
}
void vesa_swap_buffers() {
    uint32_t total_pixels = fb_width * fb_height;
    uint32_t* src = backbuffer;
    uint32_t* dst = framebuffer;
    
    // Hardware-accelerated memory copy
    asm volatile (
        "cld; rep movsl"
        : "+S" (src), "+D" (dst), "+c" (total_pixels)
        :
        : "memory"
    );
}


void put_pixel(uint32_t x, uint32_t y, uint32_t color) {
    if (x >= fb_width || y >= fb_height) return;
    uint32_t index = y * fb_width + x;
    backbuffer[index] = color;
    // uint32_t offset = y * (fb_pitch / 4) + x;
    // framebuffer[offset] = color;
}
uint8_t color_r(uint32_t c){
    return (c >> 16) & 0xFF;
}
uint8_t color_g(uint32_t c){
    return (c >> 8) & 0xFF;
}
uint8_t color_b(uint32_t c){
    return c & 0xFF;
}
uint32_t get_pixel(uint32_t x, uint32_t y) {
    if (x >= fb_width || y >= fb_height) return 0;
    uint32_t offset = y * (fb_pitch / 4) + x;
    return framebuffer[offset];
}

void draw_gradient_rect(uint32_t x, uint32_t y, uint32_t w, uint32_t h, uint32_t color1, uint32_t color2){
    for(uint32_t i = 0; i < h; i++){
        uint8_t r1 = color_r(color1);
        uint8_t r2 = color_r(color2);

        uint8_t g1 = color_g(color1);
        uint8_t g2 = color_g(color2);

        uint8_t b1 = color_b(color1);
        uint8_t b2 = color_b(color2);

        uint8_t r = (int)r1 + (((int)r2 - (int)r1) * (int)i) / (int)h;
        uint8_t g = (int)g1 + (((int)g2 - (int)g1) * (int)i) / (int)h;
        uint8_t b = (int)b1 + (((int)b2 - (int)b1) * (int)i) / (int)h;

        uint32_t color = (r << 16) | (g << 8) | b;
        draw_rect(x, y + i, w, 1, color);
    }
}
void draw_rect(uint32_t x, uint32_t y, uint32_t width, uint32_t height, uint32_t color){
    // Safety clipping
    if (x >= fb_width || y >= fb_height) return;
    if (x + width > fb_width) width = fb_width - x;
    if (y + height > fb_height) height = fb_height - y;
    
    for (uint32_t i = 0; i < height; i++) {
        // Calculate the row exactly once
        uint32_t* row_ptr = &backbuffer[(y + i) * fb_width + x];
        
        // Fast fill across the row
        for (uint32_t j = 0; j < width; j++) {
            row_ptr[j] = color;
        }
    }
}

//Bresenham's line algo...
void draw_line(int x0, int y0, int x1, int y1, uint32_t color) {
    int dx = (x1 > x0) ? (x1 - x0) : (x0 - x1); //abs(x1 - x0);
    int sx = x0 < x1 ? 1 : -1;
    int dy = (y1 > y0) ? (y0 - y1) : (y1 - y0);//-abs(y1 - y0);
    int sy = y0 < y1 ? 1 : -1;
    int err = dx + dy;
    int e2;

    while (1) {
        put_pixel(x0, y0, color);
        
        if (x0 == x1 && y0 == y1) {
            break;
        }
        
        e2 = 2 * err;
        
        if (e2 >= dy) {
            err += dy;
            x0 += sx;
        }
        if (e2 <= dx) {
            err += dx;
            y0 += sy;
        }
    }
}
void draw_circle(int xc, int yc, int radius, uint32_t color) {
    int x = radius;
    int y = 0;
    int err = 0;

    while (x >= y) {
        put_pixel(xc + x, yc + y, color); // Octant 1
        put_pixel(xc + y, yc + x, color); // Octant 2
        put_pixel(xc - y, yc + x, color); // Octant 3
        put_pixel(xc - x, yc + y, color); // Octant 4
        put_pixel(xc - x, yc - y, color); // Octant 5
        put_pixel(xc - y, yc - x, color); // Octant 6
        put_pixel(xc + y, yc - x, color); // Octant 7
        put_pixel(xc + x, yc - y, color); // Octant 8

        if (err <= 0) {
            y += 1;
            err += 2 * y + 1;
        }
        
        if (err > 0) {
            x -= 1;
            err -= 2 * x + 1;
        }
    }
}
void fill_circle(int xc, int yc, int radius, uint32_t color) {
    int x = radius;
    int y = 0;
    int err = 0;

    while (x >= y) {
        // Draw horizontal lines connecting the mirrored points
        draw_line(xc - x, yc + y, xc + x, yc + y, color);
        draw_line(xc - x, yc - y, xc + x, yc - y, color);
        draw_line(xc - y, yc + x, xc + y, yc + x, color);
        draw_line(xc - y, yc - x, xc + y, yc - x, color);

        if (err <= 0) {
            y += 1;
            err += 2 * y + 1;
        }
        
        if (err > 0) {
            x -= 1;
            err -= 2 * x + 1;
        }
    }
}
/*
void draw_char(char c, uint32_t x, uint32_t y, uint32_t fg_color, uint32_t bg_color){
    for(uint8_t row = 0; row < 8; row++){
        for(uint8_t col = 0; col < 8; col++){
            uint8_t row_data = font8x8_basic[(uint8_t)c][row];
            if(row_data >> col){
                put_pixel(x + col, y + row, fg_color);
            }else{
                put_pixel(x + col, y + row, bg_color);
            }
        }
    }
}
void draw_string(const char* str, uint32_t x, uint32_t y, uint32_t fg_color, uint32_t bg_color){
    uint32_t start_x = x;
    while (*str != '\0') {
        if (*str == '\n') {
            x = start_x;
            y += 8;
        } else {
            draw_char(*str, x, y, fg_color, bg_color);
            x += 8;   
        }
        str++;
    }
}
*/

void draw_char(char c, uint32_t x, uint32_t y, uint32_t fg_color, [[maybe_unused]]uint32_t bg_color, int scale) {
    uint8_t* glyph = font8x8_basic[(uint8_t)c];
    
    for (int row = 0; row < 8; row++) {
        for (int col = 0; col < 8; col++) {
            if ((glyph[row] >> col) & 1) {
                draw_rect(x + (col * scale), y + (row * scale), scale, scale, fg_color);
            } else {
                // draw_rect(x + (col * scale), y + (row * scale), scale, scale, bg_color);
            }
        }
    }
}

void draw_string(const char* str, uint32_t x, uint32_t y, uint32_t fg_color, uint32_t bg_color, int scale){
    uint32_t start_x = x;
    while(*str != '\0'){
        if (*str == '\n') {
            x = start_x;
            y += (8 * scale); 
        } else {
            draw_char(*str, x, y, fg_color, bg_color, scale);
            x += (8 * scale);
        }
        str++;
    }
}



\end{lstlisting}

\subsection{vesa.h}
\begin{lstlisting}[language={C++}]
#pragma once
#include <stdint.h>

uint8_t color_r(uint32_t c);
uint8_t color_g(uint32_t c);
uint8_t color_b(uint32_t c);
uint32_t get_pixel(uint32_t x, uint32_t y);
extern uint32_t* framebuffer;
extern uint32_t* backbuffer; 

// void vesa_init(uint32_t fb_addr, uint32_t width, uint32_t height, uint32_t pitch);
void vesa_swap_buffers();
void draw_gradient_rect(uint32_t x, uint32_t y, uint32_t w, uint32_t h,uint32_t color1, uint32_t color2);
void vesa_init(uint32_t fb_addr, uint32_t width, uint32_t height, uint32_t pitch);
void put_pixel(uint32_t x, uint32_t y, uint32_t color);
void draw_rect(uint32_t x, uint32_t y, uint32_t width, uint32_t height, uint32_t color);
void draw_line(int x0, int y0, int x1, int y1, uint32_t color);
void draw_circle(int xc, int yc, int r, uint32_t color);
void fill_circle(int xc, int yc, int r, uint32_t color);
void draw_char(char c, uint32_t x, uint32_t y, uint32_t fg_color, uint32_t bg_color, int scale);
void draw_string(const char* str, uint32_t x, uint32_t y, uint32_t fg_color, uint32_t bg_color, int scale);



\end{lstlisting}

\subsection{vfs.cpp}
\begin{lstlisting}[language={C++}]
#include "vfs.h"

uint32_t vfs_read(fs_node* node, uint32_t offset, uint32_t size, char* buf){
    if(node->read != 0){
        return node->read(node, offset, size, buf);
    }
    return 0;
}
\end{lstlisting}

\subsection{vfs.h}
\begin{lstlisting}[language={C++}]
#pragma once
#include <stdint.h>
#define VFS_FILE 0x01
#define VFS_DIRECTORY 0x02

struct fs_node{
    char name[128];
    uint32_t flags; // like is the file a file or a directory
    uint32_t length;
    uint32_t inode;
    uint32_t (*read)(fs_node*, uint32_t offset, uint32_t size, char* buf);
};
uint32_t vfs_read(fs_node* node, uint32_t offset, uint32_t size, char* buf);
\end{lstlisting}

\subsection{vga.cpp}
\begin{lstlisting}[language={C++}]
#include "vga.h"
#include "vesa.h"
// volatile unsigned short* vga = (volatile unsigned short*)0xB8000;
// int row = 0, col = 0;
// int cur_pos = 0;
uint32_t TERM_START_X = 0;
uint32_t TERM_START_Y = 0;
uint32_t TERM_MAX_X = 1024;
uint32_t TERM_MAX_Y = 768;

uint32_t term_x = TERM_START_X;
uint32_t term_y = TERM_START_Y;

const int FONT_SCALE = 2;
const uint32_t TERM_FG = 0xFFFFFF;
const uint32_t TERM_BG = 0x000000;


void vga_clear(){
    // for(int i = 0; i<2000; i++){
    //     vga[i] = (0x0F << 8) | ' ';
    // }
    // row = 0;
    // col = 0;
    // cur_pos = 0;
    draw_rect(TERM_START_X, TERM_START_Y, TERM_MAX_X, TERM_MAX_Y, TERM_BG);
    term_x = TERM_START_X;
    term_y = TERM_START_Y;
}
// void vga_putchar(char ch){
//     if(col == 80){
//         col = 0;
//         row++;
//         vga_putchar(ch);
//         return;
//     }
//     if(ch == '\n'){
//         col = 0;
//         row++;
//     }else{
//         vga[row * 80 + col] = (0x0F << 8) | ch;
//         col++;
//     }
// }
void scroll(){
    // if(cur_pos >= 2000){
    //     for(int i = 0; i < 24 * 80; i++){
    //         vga[i] = vga[i + 80];
    //     }
    //     for (int i = 24 * 80; i < 25 * 80; i++) {
    //         vga[i] = (0x0F << 8) | ' '; 
    //     }
    //     cur_pos = 24 * 80;
    // }
    if(term_y >= TERM_MAX_Y) vga_clear();
}
void vga_print_char(const char c){
    if(term_x >= TERM_MAX_X){
        // vga_print_char('\n');
        term_x = TERM_START_X;
        term_y += 8 * FONT_SCALE;
    }
    if(c == '\b'){
        // if(cur_pos > 0){
        //     cur_pos--;
        //     vga[cur_pos] = (0x0F << 8) | ' ';
        // }
        if(term_x > TERM_START_X){
            term_x -= 8 * FONT_SCALE;
            draw_rect(term_x, term_y, 8 * FONT_SCALE, 8 * FONT_SCALE, TERM_BG);
        }
    } else if (c == '\n') {
        // cur_pos = cur_pos + 80 - (cur_pos % 80);
        term_x = TERM_START_X;
        term_y += 8 * FONT_SCALE;
    } else {
        // vga[cur_pos] = (0x0F << 8) | c;
        // cur_pos++;
        draw_char(c, term_x, term_y, TERM_FG, TERM_BG, FONT_SCALE);
        term_x+=8 * FONT_SCALE;
    }
    scroll();
}
void vga_print(const char* s){
    for(int i = 0; s[i] != '\0'; i++){
        vga_print_char(s[i]);
    }
}
\end{lstlisting}

\subsection{vga.h}
\begin{lstlisting}[language={C++}]
#pragma once
void vga_clear();
void vga_putchar(char);
void vga_print_char(const char);
void vga_print(const char*);
\end{lstlisting}

\subsection{linker.ld}
\begin{lstlisting}[language={bash}]
ENTRY(_start)

SECTIONS{
    . = 0x100000;
    .text : { *(.multiboot) *(.text) }
    .rodata : { *(.rodata) }
    .data : { *(.data) }
    .bss : { *(COMMON) *(.bss) }
}
\end{lstlisting}


\section{Phase 26: Interactive File Reading}
In Phase 26, we transformed the static GUI into an interactive File Explorer. 

\subsection{Root Directory Chain Traversal}
Initially, the FAT32 driver only read the first sector of the Root Directory, limiting the system to 16 file entries (or fewer if Long File Names were present). This was resolved by implementing a standard FAT cluster chain traversal using a \texttt{while} loop that continuously queried the File Allocation Table for the next cluster until the EOF marker (\texttt{0x0FFFFFF8}) was reached. This allowed the system to parse massively fragmented and heavily populated root directories across multiple physical disk sectors.

\subsection{Interactive GUI Integration}
We integrated mouse state tracking (comparing current vs. previous frame \texttt{mouse\_left\_pressed} states) to detect localized clicks within the \texttt{ui\_render} File Explorer bounds. Upon detecting a collision between the mouse cursor and a file's bounding box, the GUI triggered a physical ATA sector read from the virtual \texttt{disk.img} into a shared static memory buffer. 

\subsection{The Notepad Application}
A rudimentary Window Manager visibility system was implemented by adding a boolean \texttt{visible} flag to the core \texttt{Window} structure. When a file was clicked, the hidden Notepad window was toggled to visible and popped to the front of the rendering stack. The Notepad application successfully read the shared memory buffer and utilized the existing \texttt{draw\_string} routine to rasterize the physical file contents directly onto the graphical desktop.

\end{document}
